\documentclass[preprint,12pt]{elsarticle}
\usepackage{lmodern}
\input{glyphtounicode}
\pdfgentounicode=1
\usepackage{amsmath,amssymb,amsfonts,mathtools,amsthm,bm}
\usepackage{booktabs,graphicx,array,tabularx,placeins,float,url}
\usepackage[hidelinks]{hyperref}
\hypersetup{pdfauthor={Ali Eslami and Jiangbo Yu}}
\usepackage[protrusion=true,expansion=false]{microtype}
\graphicspath{{figures/}{./}}
\allowdisplaybreaks
\setlength{\emergencystretch}{3em}
\interdisplaylinepenalty=2500
\newtheorem{assumption}{Assumption}
\newtheorem{theorem}{Theorem}
\newtheorem{proposition}{Proposition}
\newtheorem{corollary}{Corollary}
\newtheorem{remark}{Remark}
\newtheorem{definition}{Definition}
\newcommand{\R}{\mathbb{R}}
% elsarticle includes the word ``Appendix'' in appendix-section references.
\newcommand{\appref}[1]{\ref{#1}}
\newcommand{\longpaperdetails}[1]{#1 of the \href{https://github.com/AlienEslami/agentic-perimeter-runtime-certification/blob/main/paper/main.pdf}{long version}}
\newcommand{\reproductiondetails}[1]{\href{https://github.com/AlienEslami/agentic-perimeter-runtime-certification/blob/main/paper/supporting_materials/REPRODUCING_THE_REDUCED_PAPER.md}{#1}}
\journal{Transportation Research Part B: Methodological}



\begin{document}
\begin{frontmatter}
\title{Formal Guarantees for Trustworthy Agent-Mediated Urban Perimeter Control}
\author[mcgill]{Ali Eslami}
\ead{ali.eslami@mcgill.ca}
\author[mcgill]{Jiangbo Yu\corref{cor1}}
\ead{jiangbo.yu@mcgill.ca}
\cortext[cor1]{Corresponding author.}
\affiliation[mcgill]{organization={Department of Civil Engineering, McGill University},
 addressline={817 Sherbrooke Street West}, city={Montreal}, state={Quebec},
 postcode={H3A 0G4}, country={Canada}}
\begin{abstract}\fontsize{11}{13}\selectfont
Urban perimeter control must manage regional congestion and boundary queues
while satisfying time-specific service obligations. AI agents can support
traffic operators by interpreting evolving network conditions, incidents,
and service priorities and translating them into adaptive control and
replanning requests. However, each new agent-proposed plan must remain
compatible with the evolving traffic state, uncertainty, and outstanding
service obligations. We therefore formulate a runtime certification problem
for agent-proposed perimeter-control plans and develop a
transportation-specific framework that provides formal guarantees of
physical and service feasibility before execution. The framework checks
whether a proposed plan respects vehicle conservation, nonnegative traffic
states, boundary-storage limits, and service deadlines under uncertainty. It
then decides whether the plan can be executed as proposed, should be
modified to restore feasibility, or should be rejected while the previously
accepted plan remains in use, and, when the terminal conditions are
certified, ensures that accepted plans end in a state from which feasible
operation can continue. Numerical studies show that the certificate leaves
feasible plans unchanged, traffic-specific structure reduces conservatism,
and checked repair recovers feasible alternatives. Further experiments show
that the framework can accept beneficial changes while rejecting requests
that would violate boundary-storage or service constraints, including under
computation delays and agent-context errors. Overall, the framework advances
trustworthy AI for transportation systems by providing formal guarantees
that connect adaptive agent decisions to physically and operationally
feasible traffic control.
\end{abstract}
\begin{keyword}
Agentic AI \sep Macroscopic fundamental diagram \sep Perimeter control \sep Runtime certification \sep Model predictive control
\end{keyword}
\end{frontmatter}


\section{Introduction}
\label{sec:new-introduction}

Urban perimeter control regulates flows into and between regions as
congestion and external arrivals change. Regional accumulation affects
network traffic production through the macroscopic fundamental diagram
(MFD), which provides a tractable network-scale description of congestion
and trip completion
\cite{daganzo2007gridlock,geroliminis2008existence,daganzo2008analytical,geroliminis2011properties,leclercq2014cross}.
From an operator's perspective, the problem is to regulate regional
accumulation, manage vehicles waiting at region boundaries, and satisfy
time-specific service obligations. These requirements become more difficult
to coordinate when incidents, demand conditions, or operating priorities
change during execution.

Classical MFD perimeter-control methods address regional congestion through
feedback gating and multivariable boundary-flow control
\cite{keyvan2012gating,aboudolas2013perimeter}. Stability analysis and
control-Lyapunov-function designs establish regulation properties under
prescribed traffic dynamics \cite{haddad2012stability,zhong2018robust},
while robust, hybrid, and adaptive formulations extend perimeter regulation
to uncertain observations, switching actions, and changing operating
conditions
\cite{mohajerpoor2020hinf,hajiahmadi2015hybrid,haddad2020adaptive,kouvelas2017adaptive}.

Predictive perimeter control optimizes sequences of gate commands over an
MFD model \cite{geroliminis2013mpc}. Multi-region formulations account for
spatial heterogeneity and route guidance
\cite{ramezani2015heterogeneity,sirmatel2018economic}, while exterior queues
and cordon storage capture the waiting cost and spillback created by
metering \cite{haddad2017queues,ni2019cordon}. State and demand estimation
further support predictive operation when traffic conditions must be
reconstructed from measurements \cite{sirmatel2019estimation}, and recent
MFD calibration work explicitly characterizes loading--recovery variation
and uncertainty \cite{ma2025mfduq}.

AI-based decision tools provide an additional way to adapt perimeter control
as traffic conditions and operator objectives change. Data-driven
perimeter-control methods have used adaptive dynamic programming and
reinforcement learning to learn gating policies from traffic experience
\cite{chen2022dataefficient,chen2024iadp}, while multi-agent approaches
coordinate regional metering, heterogeneous cordon gates, and decentralized
control \cite{zhou2023scalable,yu2025physicsperimeter,li2024beyond}. Related
traffic-signal and language-model studies show how AI systems can generate
actions, interpret contextual information, coordinate decisions, and invoke
external tools
\cite{wei2019presslight,wei2019colight,llmlight2023,masri2025llmtraffic,collmlight2025,trafficgpt2024,openti2024,virtualpolice2026,yu2025statespace}.

These capabilities can support traffic operators when changing network
conditions or service priorities require new control decisions. In this
work, an AI agent serves this higher-level role: it interprets evolving
operating priorities and invokes computational tools, such as numerical
planners, to generate candidate perimeter-control plans.

This added flexibility creates a certification requirement whenever a
proposed plan replaces the plan currently in operation. Each agent-proposed
plan must remain compatible with the traffic state at the time of
activation, uncertainty in that state and future arrivals, available
boundary storage, and any service obligations that are still active.
Evaluation of a newly proposed plan therefore incorporates the traffic state
and outstanding obligations inherited from the operation already under way.
The operator needs a principled way to decide whether the new plan may
replace the current one, whether it should be modified, or whether the
previously accepted plan should remain in operation.

Runtime assurance provides a natural foundation for this decision. Simplex
architectures retain a verified fallback controller
\cite{simplex,mehmood2022black}; control barrier functions and shielding
restrict unsafe actions \cite{ames2017,alshiekh2018}; predictive safety
filters evaluate future constraint satisfaction
\cite{wabersich2021,hsu2023safety}; and runtime-assurance and
reference-governor methods formalize the transfer of control authority
\cite{fisac2019,hobbs2023runtime,gilbert2002generalized,garone2017reference}.
Reachability methods additionally provide tools for propagating uncertainty
over future trajectories \cite{coogan2020mixed,kochdumper2021sparse}. These
methods establish foundations for supervising learning-enabled decisions.
Urban perimeter control also provides structure that can be exploited
directly: vehicles wait at region boundaries for admission to particular
destinations, regional inventories and boundary queues are physically
coupled, and service requirements may remain active across successive
planning requests.

We use this transportation structure to construct a runtime certificate for
agent-proposed perimeter-control plans. The certificate checks whether a
candidate plan satisfies vehicle conservation, nonnegative traffic states,
boundary-storage limits, and time-specific service requirements for all
states and arrivals represented by the uncertainty model. A key observation
is that a vehicle waiting for admission to a destination and the
corresponding destination inventory form a conserved quantity with respect
to that admission flow. Combining these quantities therefore removes the
admission decision from part of the uncertainty propagation. Positivity,
independent queue bounds, and signed uncertainty representations further
tighten the resulting bounds while preserving dependencies in the uncertain
traffic state. Together, these ingredients provide a transportation-specific
representation for checking regional inventories, queues, control inputs,
and service constraints over the planning horizon.

The certificate is then integrated with sequential planning and execution.
If a proposed plan satisfies the required conditions, it can replace the
currently active plan with its commands preserved. A rejected plan enters a
checked repair procedure that searches for a nearby feasible alternative,
with every repaired candidate returned to the same certificate before
execution. If no replacement is accepted, the previously accepted plan
remains active through its certified validity interval. When the terminal
conditions are certified, the framework ensures that accepted plans end in a
state from which feasible operation can continue. It also accounts for
traffic evolution while a new plan is being computed and certified.

The resulting framework provides a formal interface between adaptive
AI-based planning and physical traffic execution. The interface determines
when plans produced by AI models and computational tools are admissible for
the transportation system. The numerical evaluation examines certification
conservatism and checked repair. Further experiments show that the framework
can accept beneficial changes while rejecting requests that would violate
boundary-storage or service constraints, including under computation delays
and agent-context errors.

The paper makes four main contributions:
\begin{enumerate}
\item We define the requirements that agent-proposed perimeter-control plans
must satisfy before execution. These requirements connect operator requests
to numerical planning tools and ensure that each new plan accounts for the
current traffic state, uncertainty, existing service obligations, and plan
timing.

\item We develop a traffic-specific method for checking physical constraints
and service deadlines under uncertainty. The method uses vehicle
conservation, nonnegative vehicle counts, and the links between boundary
queues and regional traffic to obtain tighter bounds on future traffic
states. It retains relationships between uncertain traffic quantities and
extends to networks with several entry queues.

\item We prove guarantees for accepting, repairing, and rejecting proposed
plans. Plans that pass certification can be executed unchanged. With a
backup certified for the same traffic information and requirements, a
completed repair search selects the certified plan closest to the proposed
plan within a predefined finite set, measuring the difference through their
planned control actions. If a proposal is rejected or computation runs out
of time, the previously accepted plan remains in use through its certified
validity interval. When the required terminal conditions are also certified,
feasible operation can continue beyond plan expiry.

\item We provide reusable test cases involving regional congestion, boundary
queues, and service deadlines. Numerical studies examine the tightness of
the certification bounds and how well repair recovers feasible plans.
Further experiments test successive plan changes, networks with several
regions, computation delays, and errors in the information given to the
agent.
\end{enumerate}

Section~\ref{sec:new-problem} formulates the traffic-control problem. Section~\ref{sec:new-agent} defines the agent and planning interface, and Section~\ref{sec:new-certification} develops plan certification and execution. Section~\ref{sec:new-evidence} evaluates the method; Section~\ref{sec:new-discussion} discusses operational implications. Complete proofs accompany the results in Section~\ref{sec:new-certification}; essential experimental settings appear in the appendix, and detailed reproduction material is available in the repository.

\begin{table}[!tp]
\centering\small
\caption{Principal notation. Traffic inventories are expressed in units of
100 vehicles unless a table states physical vehicles.}
\label{tab:nomenclature}
\renewcommand{\arraystretch}{0.98}
\begin{tabular}{@{}p{0.20\linewidth}p{0.72\linewidth}@{}}
\toprule
Symbol & Meaning\\
\midrule
$n_i$, $n_{ij}$ & Regional accumulation; vehicles in region $i$ destined for $j$\\
$z$, $x$, $M$ & Exterior waiting inventory; physical state; total vehicle inventory\\
$u_{ij}$, $a$ & Fractional gate command; realized exterior admission flow\\
$\Lambda$, $d^{\rm int}_{ij}$ & Exterior arrivals; destination-specific internal demand\\
$d$, $\bar w_{\rm veh}$ & Sensing half-width (vehicles); physical arrival-residual bound (veh/h)\\
$\bar w_x$ & Normalized arrival-residual bound $\bar w_{\rm veh}/100$ (inventory units/h)\\
$\mathcal S_t$ & Decision record: state set, uncertainty, task, active plan, time and versions\\
$P$, $P^B$, $\mu_B$ & Proposed plan; verified continuation; terminal backup policy\\
$P_a$, $c$ & Destination-allocation matrix; network destination-class vector\\
$t_f$, $\eta$ & Example service deadline and inventory limit\\
$T_r$, $g_r$, $\eta_r$ & General service time, row vector and threshold\\
$\theta_i$, $\lambda_{\rm rep}$ & Agent task record; predictive-repair interpolation weight\\
$K_b$, $K$, $\mathcal K$ & Feedback gain; local block count; terminal invariant set\\
$\beta_j^{\rm phys}$, $b_a$ & Physical constraint-row limit; admission-incidence vector\\
$y$, $L$ & Conservation-coordinate state and its invertible transform\\
$r_l$, $R_l$ & Starting and trial-tube radii on interval $l$\\
$a_{il}$, $b_{ijl}$ & Local diagonal and off-diagonal Jacobian bounds\\
$\xi$, $\Gamma_l$, $\zeta_l$ & Initial uncertainty, signed generator and nonlinear remainder\\
$z_e$ & Prediction-error coordinate used only in the enclosure proof\\
$J_R$, $J_Q$, $J_D$ & Native regulation, queue-economic and weighted-priority objectives\\
$p$, $E$, $x^R$ & Regional-priority weight; finite-plan endpoint; native four-state vector\\
$\mathcal C$, $\mathcal F_{\mathcal C}$, $\widehat{\mathcal F}_{\mathcal C}$ & Planning contract; robust feasible and checker-certified plan sets\\
$X_i^A$, $h_i$, $\iota_i$ & Agent interface state, memory and supplied contextual view\\
$B_\omega$ & Fraction of requested priority benefit recovered by repair\\
$\mathrm{TTS}$ & Total time spent, including exterior waiting; vehicle-hours (veh-h)\\
\bottomrule
\end{tabular}
\end{table}

\section{Two-Region MFD Perimeter-Control Problem}
\label{sec:new-problem}

The two-region model describes destination-specific traffic, interregional
transfers and an exterior entry queue. Section~\ref{sec:network-extension}
extends the construction to several regions and queues.

\subsection{Traffic inventories, flows and boundary storage}
\label{app:new-model}

The state and gate commands are
\begin{equation}
\begin{aligned}
x&=(n_1,n_2,n_{11},n_{22},n_{10},n_{20},z)^\top,\\
u&=(u_{12},u_{21},u_{02},u_{20})^\top.
\end{aligned}
\label{eq:new-state}
\end{equation}
Here $n_i$ is the total inventory in region $i$, and $n_{ij}$ counts vehicles
in region $i$ destined for $j$. Destination $0$ denotes the exterior, and
$n_{ii}$ counts vehicles currently in region $i$ whose trips will end there.
Exterior-destined vehicles in
Region~1 leave through Region~2. The remaining classes are
$n_{12}=n_1-n_{11}-n_{10}$ and $n_{21}=n_2-n_{22}-n_{20}$.
The queue $z$ holds vehicles awaiting admission. Inventories are continuous
aggregate quantities in units of 100 vehicles, with time measured in hours.

Following \cite{zhong2018robust,geroliminis2013mpc}, potential outflow is
\begin{equation}
\begin{aligned}
G(n)&=0.0014877n^3-0.29815n^2+15.0912n,\\
q_{ij}&=n_{ij}v(n_i),
\end{aligned}
\label{eq:new-mfd}
\end{equation}
where $v(n)=0.0014877n^2-0.29815n+15.0912$ extends $G(n)/n$ continuously
to zero. Thus $n_{ij}$ is an inventory and $q_{ij}$ is a potential flow
before gating, in inventory units/h. The MFD allocates potential trip
completion and regional transfer among destination classes in proportion
to their inventories. Its critical accumulation is 3391.93 vehicles and
peak outflow is 22,691 veh/h. The operational cap of 10,000 vehicles per
region lies within the fitted MFD's positive-production range. The model
assumes comparable within-region trip lengths, destination-proportional
mixing and adequate receiving supply on the verified domain.

Exterior arrivals are $\Lambda(t)=\widehat\Lambda(t)+w(t)$, with forecast
$\widehat\Lambda$ and bounded measurable residual $w$. A fraction
$\pi=1/2$ is destined for Region~1. Admission and queue dynamics are
\begin{equation}
a(z,u_{02})=\min\{60z,120u_{02}\},\qquad
\dot z=\Lambda-a.
\label{eq:new-admission}
\end{equation}
The two branches represent a one-minute release time and commanded
admission capacity. The queue models aggregate boundary storage
\cite{haddad2017queues}. Spatial spillback requires receiving-supply dynamics.
Gate commands are dimensionless fractions: $u_{ij}=1$ permits the full
potential class flow, while $u_{ij}=0$ blocks it. Exterior gate $u_{02}$
instead scales the admission capacity. Commands vary over time or are held
on finite blocks. For a class flow $q$ and gate $u_g$, the number transferred
over $[t_a,t_b]$ is $100\int_{t_a}^{t_b}u_g(t)q(t)\,dt$.

Let $d^{\rm int}_1=d^{\rm int}_{11}+d^{\rm int}_{12}+d^{\rm int}_{10}$ and $d^{\rm int}_2=d^{\rm int}_{22}+d^{\rm int}_{21}+d^{\rm int}_{20}$ be fixed
internal demands. Vehicle conservation gives
\begin{equation}
\begin{aligned}
\dot n_1&=-q_{11}-(q_{12}+q_{10})u_{12}+q_{21}u_{21}+d^{\rm int}_1,\\
\dot n_2&=-q_{22}+(q_{12}+q_{10})u_{12}-q_{21}u_{21}
                 +a-q_{20}u_{20}+d^{\rm int}_2,\\
\dot n_{11}&=-q_{11}+q_{21}u_{21}+d^{\rm int}_{11},\\
\dot n_{22}&=-q_{22}+q_{12}u_{12}+(1-\pi)a+d^{\rm int}_{22},\\
\dot n_{10}&=-q_{10}u_{12}+d^{\rm int}_{10},\\
\dot n_{20}&=-q_{20}u_{20}+q_{10}u_{12}+d^{\rm int}_{20},\\
\dot z&=\widehat\Lambda+w-a.
\end{aligned}
\label{eq:new-plant}
\end{equation}
A transfer preserves destination: vehicles leaving $n_{21}$ enter
$n_{11}$, already included in $n_1$, and later complete their trips through
$q_{11}$. Inventories remain continuous when commands change. Commands obey
\begin{equation}
\mathcal U=\{u:0\le u\le1,\ |u_{21}-u_{12}|\le0.75,
\ |u_{20}-u_{02}|\le0.75\}.
\label{eq:new-gates}
\end{equation}
Nonnegative classes, $0\le n_1,n_2\le100$, $0\le z\le20$ and
\eqref{eq:new-gates} give 24 affine state--command inequalities.

The initial state is $(30,30,15,15,5,5,2)$ and base exterior arrivals are
6000 veh/h. In the order $(11,12,10,22,21,20)$, fixed internal demands are
\begin{equation}
\begin{aligned}
\boldsymbol d^{\rm int}=(&74.85630,37.42815,18.714075,\\
              &44.85630,7.42815,0.748563).
\end{aligned}
\label{eq:new-incident-demands}
\end{equation}
They are constructed from this equilibrium with
$u^\star=(0.5,0.5,0.5,0.52)$ and remain fixed across events and reference
choices. Displayed decimals are rounded. Each proof uses the stored
coefficients of the implemented field. Reference-controller construction
and exact coefficients are supplied in the \reproductiondetails{model records}.

\subsection{Service requirements and planning tasks}
\label{sec:benchmark-formulation}

Service requirements are scheduled affine rows $g^\top x(t_f)\le\eta$.
For example, a representative incident requires
\begin{equation}
100n_1(t_f)\le2850\ \text{vehicles},\qquad t_f=1\ \text{h}.
\label{eq:new-service}
\end{equation}
The binding-service study uses 1700 vehicles at the same deadline. Only
the operator can change a service target or its deadline. Changing a soft
planning objective preserves these requirements.
Adding the regional and queue balances gives
\begin{equation}
\begin{aligned}
M&=n_1+n_2+z,\\
\dot M&=d^{\rm int}_1+d^{\rm int}_2+\Lambda-F_{\rm out},\\
F_{\rm out}&=q_{11}+q_{22}+q_{20}u_{20},\\
\mathrm{TTS}&=100\int_0^{T_{\rm eval}} M(t)\,dt.
\end{aligned}
\label{eq:new-conservation}
\end{equation}
Total time spent (TTS) therefore includes exterior waiting. Completed
departures are $100\int F_{\rm out}\,dt$. For matched arrivals and initial
inventory, their difference is $-100\Delta M(T_{\rm eval})$.

\emph{Task Q} requests two constant command blocks from $t_0=1/3$ h to
$E=1$ h, subject to the physical rows, minimizing
\begin{equation}
J_Q(P)=100\int_{t_0}^{E}(n_1+n_2+z)\,dt.
\label{eq:benchmark-queue-cost}
\end{equation}
\emph{Task D} adds $100n_1(E)\le\eta$ and a soft Region-2 priority with
weight $p\ge1$, minimizing
\begin{equation}
J_D(P)=100\int_{t_0}^{E}(n_1+p n_2+z)\,dt.
\label{eq:benchmark-priority-cost}
\end{equation}
For sensing half-width $d$ vehicles and arrival-residual bound
$\bar w_{\rm veh}$ veh/h, let $\bar w_x=\bar w_{\rm veh}/100$ and use
\begin{equation}
\begin{aligned}
\mathcal I_{t_0}&=\widehat x_{t_0}+[-d/100,d/100]^7,\\
|w(t)|&\le\bar w_x,\qquad t\in[t_0,E].
\end{aligned}
\label{eq:benchmark-contract}
\end{equation}
The supplementary regional-regulation task R, including its native
four-state model and squared-error objective, is specified in the
\reproductiondetails{additional controller comparison}.

A contract $\mathcal C$ fixes the state set, uncertainty family, model,
physical rows, service rows and interval $[t_0,E]$. For the executable plans
of Section~\ref{sec:new-agent}, write $x_P(s;x_0,w)$ for the trajectory.
The robust feasible set is
\begin{equation}
\begin{aligned}
\mathcal F_{\mathcal C}=\{P:\;&\forall x_0\in\mathcal I_{t_0},\ w\in\mathcal W,\\
&H_jx_P(s;x_0,w)+J_ju_P(s;x_0,w)
                         \le\beta_j^{\rm phys}\\
&\hspace{1em}\text{for all }s\in[t_0,E]\text{ and physical rows }j;\\
&g_r^\top x_P(T_r;x_0,w)\le\eta_r
\quad\text{for all service rows }r\}.
\end{aligned}
\label{eq:benchmark-feasible-set}
\end{equation}
Membership requires a well-defined trajectory throughout $[t_0,E]$.
$\mathcal W$ contains every measurable residual satisfying the bound, and
all service times lie in this interval. Task Q has no service rows.
Task D includes the minute-60 row. The checker returns a computable subset
\begin{equation}
\widehat{\mathcal F}_{\mathcal C}
=\{P:\mathsf{Cert}_{t_0}(P\mid\mathcal C)=1\}
\subseteq\mathcal F_{\mathcal C}.
\label{eq:benchmark-certifiable-set}
\end{equation}
Section~\ref{sec:new-certification} proves this containment. Planners
minimize a center-forecast cost under nominal or finite-scenario
constraints. Certification checks the full uncertainty family. We write
$J_Q(P;x_0,w)$ or $J_D(P;x_0,w)$ when comparing particular realizations.
A sampled violation disproves feasibility, whereas passing sampled
trajectories alone leaves full-family feasibility undecided. A replacement
or reserve for a changed contract requires a certificate for that contract.

\section{Agent-Mediated Planning and Execution Authority}
\label{sec:new-agent}

The agent interprets operator instructions and requests numerical plans.
The execution layer checks each proposed plan against the current traffic
information, uncertainty and service obligations.

\subsection{Requests and task timing}
\label{app:task-state-compiler}

At proposal time $t_i$, memory $h_i$, task record $\theta_i$ and observation
and instruction $o_i$ form the interface state:
\begin{equation}
\begin{aligned}
X_i^A&=(h_i,\theta_i,o_i),\\
\iota_i&=\mathsf{View}(o_i,\theta_i),\\
r_i^A&\sim\Pi_\phi(\cdot\mid h_i,\iota_i),\\
(\theta_i^+,q_i^A)&=\mathsf{Comp}_{t_i}(\theta_i,r_i^A),\\
\mathcal P_i^{\rm prop}
&=\begin{cases}\mathsf{Planner}(q_i^A,\mathcal I_i,\mathcal T_i),&q_i^A\ne\bot,\\
\varnothing,&q_i^A=\bot,\end{cases}\\
(u_i^{\rm exec},\ell_i)
&=\mathsf{CertifyExecute}(\mathcal P_i^{\rm prop},\mathcal S_{t_i}),\\
\theta_{i+1}&=\mathsf{Age}_{t_{i+1}}(\theta_i^+),\\
h_{i+1}&\sim\mathsf U_\phi(\cdot\mid
h_i,\iota_i,r_i^A,u_i^{\rm exec},\ell_i,o_{i+1}).
\end{aligned}
\label{eq:new-agent-state}
\end{equation}
An invalid request returns $q_i^A=\bot$ and an empty proposal. The execution layer then retains its valid incumbent.
A history-only view supplies the observation and memory. A task-record
view also supplies the active priority and its absolute expiry. The model
$\Pi_\phi$ proposes an operation. Compilation checks its fields and updates
the requested soft objective, with physical information $\mathcal I_i$ and
authorized service requirements $\mathcal T_i$ supplied separately.

The task record contains an active regional priority and a last-in,
first-out stack of suspended priorities, each with region, label and
absolute expiry. Network operation is the default. Before compilation and
between proposals, $\mathsf{Age}_t$ removes an expired active task and then
discards expired predecessors until an unexpired task is restored or the
stack is empty. Aging compares absolute expiries with the current time. Enumerate the
active task and then the suspended tasks from most recent to oldest, with
expiries $\tau_1^{\mathrm{exp}},\ldots,\tau_m^{\mathrm{exp}}$. The first
surviving entry is
\begin{equation}
k_t=\min\{k\in\{1,\ldots,m\}:\tau_k^{\mathrm{exp}}>t\}.
\label{eq:task-age-survivor}
\end{equation}
If this set is empty, aging returns network operation. Otherwise entry
$k_t$ becomes active, its older predecessors remain suspended, and all
original expiries are preserved. On the active task and suspension stack,
\begin{equation}
\mathsf{Age}_t(\mathsf{Age}_s\theta)=\mathsf{Age}_t\theta,
\qquad t\ge s.
\label{eq:task-age-composition}
\end{equation}
Indeed, entries removed at $s$ have expiry at most $s\le t$ and cannot
become the first surviving entry at $t$. Setting $s=t$ proves idempotence,
so aging both between proposals and before compilation is consistent.
Audit-history timestamps are separate from this task-state identity.

For a parsed operation,
\begin{equation}
\theta_i^+=\mathsf{UpdateTask}(\mathsf{Age}_{t_i}\theta_i,r_i^A),\qquad
(j_i,\Delta_i)=\mathsf{SelectTask}(\theta_i^+,t_i).
\label{eq:task-state-compiler}
\end{equation}
Here $\mathsf{UpdateTask}$ updates the parsed task record and
$\mathsf{SelectTask}$ selects its objective and remaining duration.
The compiler $\mathsf{Comp}$ also forms the numerical planning request.
The index $j_i$ selects the objective. A task expiring at $\tau_i^{\mathrm{exp}}$ has remaining
duration $\Delta_i=\tau_i^{\mathrm{exp}}-t_i>0$, while network operation has $\Delta_i=0$.
The transitions are given in Table~\ref{tab:task-state-operations}.

\begin{table}[t]
\centering\small
\caption{Task-record transitions at current time $t$. Expiry is a
priority's absolute end time. The stack stores interrupted priorities,
most recent first, with original expiries. Network operation uses the
default network objective.}
\label{tab:task-state-operations}
\begin{tabular}{@{}p{0.16\linewidth}p{0.76\linewidth}@{}}
\toprule
Operation & Transition\\
\midrule
Set & Clear the stack and install the specified priority until $t+D/60$, for duration $D$ minutes.\\
Interrupt & Save the active task and install the temporary priority until
$t+D/60$, for duration $D$ minutes.\\
Continue & Preserve the surviving task and its expiry.\\
Extend & Add $D/60$ hours to the active task's expiry, for extension $D$ minutes.\\
Resume & Restore the latest unexpired saved task. If none survives, clear
the interruption. With no stack or interruption, keep the current task.\\
Cancel; network & Clear active and suspended tasks and use network operation.\\
\bottomrule
\end{tabular}
\end{table}

\emph{Temporal consistency.} The traffic clock $t$ is in hours.
An instruction duration $D$ or extension $D_m^{\rm ext}$ supplied in
minutes is divided by 60 before updating an expiry. Thus every stored task
has
\begin{equation}
\tau^{\rm exp}=t^{\rm create}
 +\frac{D+\sum_{m=1}^{n^{\rm ext}}D_m^{\rm ext}}{60}.
\label{eq:task-expiry-invariant}
\end{equation}
Set and interrupt establish this equality, and extend adds its stated positive
increment. Continue, suspension and restoration leave it unchanged.
Aging removes only expired entries. These observations prove the invariant
by induction over the transition table. Positive durations and selection
of unexpired predecessors give $\Delta_i>0$ for a regional task. Network
operation has $\Delta_i=0$. Invalid operations leave the aged record
unchanged and return $q_i^A=\bot$, so no planner is called. The stored
implementation uses the equivalent minute clock $60t$. These operations
change soft priorities while admitted hard service requirements remain in
$\mathcal T_i$.

A valid operation may update the requested objective while execution
retains an older certified plan. Interpretation accuracy is assessed against
the operator instruction. The certificate establishes the executed plan's
physical and service properties. The memory update $\mathsf U_\phi$ uses the
request, execution feedback and next observation. A single possible memory
state gives an agent without stored interaction history.

\subsection{Plans and execution authority}

A finite plan specifies affine command blocks $u=k_b+K_bx$ on absolute
intervals. Target-centered feedback has $k_b=u_b^\star-K_bx_b^\star$.
Constant model predictive control (MPC) commands have $K_b=0$.
MPC, feedback-rollout and repair tools
all submit their plans to the same checker (Figure~\ref{fig:new-architecture}).
The optional repair interface is
\begin{equation}
\widehat P_i^R=\mathsf{RepairTool}(P_i^0,P_i^B,\mathcal S_{t_i},\mathcal G_{\rm rep}).
\label{eq:repair-tool-interface}
\end{equation}
The decision record $\mathcal S_{t_i}$ contains the physical information,
absolute time, model and constraint versions, obligations and verified
continuation. An accepted replacement starts at its certified installation
time. Rejection preserves the current plan's unexecuted commands and original
expiry. Decision times are locally finite. Section~\ref{sec:predictive-repair}
gives the repair and installation rules. Section~\ref{sec:terminal-backup}
gives conditions for continued operation beyond expiry.

\begin{figure}[t]
\centering
\includegraphics[width=0.98\textwidth]{figures/mfd_tool_architecture.pdf}
\caption{Planning and execution. AI means artificial intelligence, MPC
means model predictive control, and MFD means macroscopic fundamental
diagram. Numerical planning tools submit candidate plans to a common
certifier. The execution authority installs an accepted plan or continues
the still-valid plan. A certified terminal policy takes over at expiry
when available. Arrows show requests, candidate plans, acceptance or
rejection decisions, executed commands and traffic feedback.}
\label{fig:new-architecture}
\end{figure}

\section{Structure-Preserving Runtime Certification}
\label{sec:new-certification}
\label{app:repair-cost}

The certificate follows vehicles inside the regions and waiting at their
boundaries. It checks one proposed plan, then preserves those guarantees
when plans are replaced, repaired, or retained. All results below apply
under the stated model, uncertainty, and execution conditions.

\subsection{Plan data and conservation coordinates}

A plan starts at $t_0$, expires at the absolute time $E$, and contains
finitely many command blocks. On block $b$,
\begin{equation}
 u_b(x)=k_b+K_bx,
 \qquad \Lambda(t)=\widehat\Lambda(t)+w(t),
 \qquad |w(t)|\le\bar w_x.
 \label{eq:new-plan-data}
\end{equation}
The primary two-region forecast is continuous and piecewise affine.
$w$ may be any measurable signal within this bound. The residual uses the
normalized flow units of Section~\ref{sec:benchmark-formulation}.
Target feedback has $k_b=u_b^\star-K_bx_b^\star$, and constant MPC commands
have $K_b=0$. Command changes preserve physical inventories.

Physical limits and scheduled service requirements are affine rows:
\begin{align}
 H_jx+J_ju&\le\beta_j^{\rm phys},\quad j=1,\ldots,n_c,
 \label{eq:new-physical-rows}\\
 g_r^\top x(T_r)&\le\eta_r,\quad r=1,\ldots,n_s.
 \label{eq:new-service-rows}
\end{align}
The physical rows include nonnegative destination classes, regional and
exterior storage, gate boxes, and paired-gate differences. On block $b$,
write $h_{jb}^\top=H_j+J_jK_b$ and
$\beta_{jb}=\beta_j^{\rm phys}-J_jk_b$, so that the row becomes
$h_{jb}^\top x\le\beta_{jb}$. With $e_i$ the $i$th coordinate vector,
the incident requirement has $g_r=100e_1$ and $\eta_r=2850$ vehicles.
Two rows represent a two-sided imbalance band. The operator fixes the
requirements, deadlines, and task version.

For the fixed destination fraction $\pi$, use
\begin{equation}
\begin{aligned}
y=Lx=(&n_{11},n_{12},n_{10},n_{22}+(1-\pi)z,\\
      &n_{21}+\pi z,n_{20},z)^\top.
\end{aligned}
\label{eq:new-conservation-coordinates}
\end{equation}
This invertible map combines each receiving class with its share of the
waiting vehicles. Its inverse includes
\begin{equation}
 \begin{aligned}
 n_1&=y_1+y_2+y_3,&n_2&=y_4+y_5+y_6-y_7,\\
 n_{11}&=y_1,&n_{22}&=y_4-(1-\pi)y_7,\\
 n_{10}&=y_3,&n_{20}&=y_6,\qquad z=y_7.
 \end{aligned}
 \label{eq:new-inverse-coordinates}
\end{equation}
For admission incidence $b_a=(0,1,0,1-\pi,0,0,-1)^\top$,
\begin{equation}
 Lb_a=-e_7,
 \qquad Le_7=(0,0,0,1-\pi,\pi,0,1)^\top.
 \label{eq:new-conservation-incidence}
\end{equation}
Adding each receiving-class balance to its share of
$\dot z=\Lambda-a$ cancels admission from the combined inventories.
Consequently $n_1+n_2+z=\sum_{i=1}^6y_i$, and explicit admission remains
only in the queue equation. Regional production still depends on the
full transformed state.

\subsection{Certification of one plan}
\label{app:new-enclosure-proof}

Let $f_b(x,t)$ be the block's field under the nominal forecast. The actual
field is $f_b(x,t)+e_7w(t)$. Class production is evaluated through the
polynomial completion factor $v(n)=G(n)/n$, extended continuously at zero.
The admission minimum has affine branches under affine feedback. Their
minimum is continuous and locally Lipschitz, so the field is locally
Lipschitz in state and measurable in time. Bounded measurable residuals
therefore give unique local absolutely continuous solutions. The compact
enclosures below continue them across each finite plan.

Split the plan at command changes, forecast knots, service times, and
prediction checkpoints. On $[t_l,t_l+h_l]$, choose a differentiable
predictor $p_l$ and inherited radii $r_l\ge0$ satisfying
$|L(x(t_l)-p_l(t_l))|\le r_l$. For an initial physical box centered at
$p(t_0)$ with half-width $d_x$, use $r_0\ge|L|d_x$. The initial predictor
box is $p(t_0)+L^{-1}[-r_0,r_0]$. A trial radius $R_l>0$ defines
$p_l(t)+L^{-1}[-R_l,R_l]$. On this whole region and time interval, obtain
\begin{equation}
 \begin{aligned}
 a_{il}&\ge\sup (L D_xf_b L^{-1})_{ii},\\
 b_{ijl}&\ge\sup |(L D_xf_b L^{-1})_{ij}|\quad(j\ne i),\\
 \delta_{il}&\ge\sup_t |(L f_b(p_l(t),t)-L\dot p_l(t))_i|.
 \end{aligned}
 \label{eq:new-local-bounds}
\end{equation}
At a branch surface the generalized Jacobian is contained in the convex
hull of limiting smooth branch Jacobians. Both branches are included
whenever the trial region meets that surface. A branch is omitted only
when verified bounds exclude it throughout the region. The diagonal bound
retains its sign, while cross-component terms use magnitudes.

The scalar comparison is
\begin{align}
 c_{il}(R_l)&=\sum_{j\ne i}b_{ijl}R_{jl}
                  +\delta_{il}+|(Le_7)_i|\bar w_x,\nonumber\\
 \phi_{il}(s)&=e^{a_{il}s}r_{il}
                   +\chi(a_{il},s)c_{il}(R_l),\nonumber\\
 \chi(a,s)&=\begin{cases}(e^{as}-1)/a,&a\ne0,\\s,&a=0.\end{cases}
 \label{eq:new-component-comparison}
\end{align}
Containment and continuous physical feasibility pass when
\begin{align}
 \max\{r_{il},\phi_{il}(h_l)\}&<R_{il}\quad\text{for every }i,
 \label{eq:new-self-inclusion}\\
 \sup_{t\in[t_l,t_l+h_l]}h_{jb}^\top p_l(t)
       +|h_{jb}^\top L^{-1}|R_l&\le\beta_{jb}.
 \label{eq:new-affine-support}
\end{align}
Absolute values are componentwise. At a boundary, bound predictor recentering by
\begin{equation}
 \delta^{\rm rec}_l\ge
 |L(p_l(t_{l+1})-p_{l+1}(t_{l+1}))|,
 \label{eq:new-recentering-bound}
\end{equation}
and propagate $r_{i,l+1}=\phi_{il}(h_l)+\delta^{\rm rec}_{il}$. At a service
checkpoint the center $\widehat x(T_r)$ and radius $r(T_r)$ give
\begin{equation}
 g_r^\top\widehat x(T_r)+|g_r^\top L^{-1}|r(T_r)\le\eta_r.
 \label{eq:new-service-support}
\end{equation}

The following theorem turns these checks into a guarantee for every
trajectory allowed by the plan's initial set and uncertainty bounds.

\begin{theorem}
\label{thm:new-plan-enclosure}
Consider the MFD field \eqref{eq:new-plant} with admission
\eqref{eq:new-admission}, a finite affine-command plan, and the uncertainty
class \eqref{eq:new-plan-data}. Let its initial information set be contained
in the initial predictor box. Suppose the bounds
\eqref{eq:new-local-bounds} hold for every state and time in each trial
region, and
\eqref{eq:new-self-inclusion} holds on every interval. Then the propagated
boxes contain every trajectory from that initial set under every allowed
residual. If \eqref{eq:new-affine-support} and
\eqref{eq:new-service-support} also pass, the trajectory satisfies every
physical row continuously through the plan's expiry and every included
service row at its deadline. At command or target changes, the inherited box
is propagated and the successor command rows are checked on that box.
\end{theorem}

\begin{proof}
Consider one interval and omit its index. Before a hypothetical first exit
from $|z_e|<R$, where $z_e=L(x-p)$, the segment between predictor and true
state lies in the trial region. The Lipschitz mean-value relation gives
\begin{equation}
 \dot z_e=A(t)z_e+\epsilon(t)+(Le_7)w(t),
 \label{eq:new-error-dynamics}
\end{equation}
almost everywhere. Here $A(t)$ is in the convex hull of transformed
generalized Jacobians along the segment, and
$\epsilon=Lf_b(p,t)-L\dot p$ satisfies $|\epsilon_i|\le\delta_i$.
The branch bounds therefore apply also to this mean-value matrix.
For $v_i=|(z_e)_i|$, they imply
\begin{equation}
 \dot v_i\le a_iv_i+c_i(R)
 \quad\text{almost everywhere before exit}.
 \label{eq:new-error-comparison-proof}
\end{equation}
At a zero component the upper derivative has the same forcing and
cross-component bound. Integrating-factor comparison gives
$v_i(t_l+s)\le\phi_i(s)$. Since
$\phi_i'(s)=e^{a_is}(a_ir_i+c_i(R))$ has constant sign, its maximum on
$[0,h]$ is $\max\{r_i,\phi_i(h)\}$. Strict
\eqref{eq:new-self-inclusion} contradicts a first exit by continuity.
The compact enclosure also prevents finite-time escape of the solution.

At the right endpoint, the old-center error is bounded by $\phi(h)$.
Physical states remain continuous at command changes. The triangle
inequality and \eqref{eq:new-recentering-bound} give the stated successor
radius. Induction across the finite partition proves containment.
For a row, write $x=p+L^{-1}z_e$. Its maximum over $|z_e|\le R$ is
$h^\top p+|h^\top L^{-1}|R$. Maximizing the predictor term in time proves
\eqref{eq:new-affine-support}. The endpoint radius proves
\eqref{eq:new-service-support}. These checks establish every continuous
physical row and every scheduled service row, including successor-command
rows at boundaries.
\end{proof}

All bounds used for acceptance include outward rounding of coefficients,
range evaluations, and comparison functions. Numerically generated
predictors and trial radii obtain authority only after these inequalities
pass. When $K_b=0$, command rows are constants: their stored binary
coefficients are treated as exact binary rationals and their residuals
checked exactly, including zero. State and storage rows still use verified
trajectory bounds. This preserves feasibility of commands on a closed gate
boundary without enlarging the admissible box.

\subsection{Traffic structure and tighter bounds}
\label{sec:structure-certificate}
\label{app:structure-proof}

Nonnegative vehicle counts, independent queue dynamics, and shared
initial uncertainty provide tighter certificates for the same constraints.
For constant commands, the derived class balances are
\begin{equation}
 \begin{aligned}
 \dot n_{12}&=d^{\rm int}_{12}-u_{12}v(n_1)n_{12},\\
 \dot n_{21}&=d^{\rm int}_{21}+\pi a-u_{21}v(n_2)n_{21}.
 \end{aligned}
 \label{eq:structure-derived-classes}
\end{equation}
For the installed polynomial coefficients $\alpha_2=0.0014877$,
$\alpha_1=-0.29815$ and $\alpha_0=15.0912$, exact arithmetic gives
$\alpha_2>0$ and $4\alpha_2\alpha_0-\alpha_1^2>0$, hence
\begin{equation}
 v(n)=\alpha_2\left(n+\frac{\alpha_1}{2\alpha_2}\right)^2+
                    \frac{4\alpha_2\alpha_0-\alpha_1^2}{4\alpha_2}>0.
 \label{eq:structure-positive-factor}
\end{equation}
This continuation of the polynomial establishes positivity independently
of the requested jam and storage limits.

The queue satisfies the independent scalar equation
\begin{equation}
 \dot z=\widehat\Lambda(t)+w(t)-\min\{60z,120u_{02}\}.
 \label{eq:structure-scalar-queue}
\end{equation}
Extreme initial queues and arrival bounds define lower and upper
comparison solutions. For constant arrival $\lambda$ and
$\kappa=120u_{02}$, each branch has the solution
\begin{equation}
 \begin{aligned}
 z(s)&=\lambda/60+(z_0-\lambda/60)e^{-60s}&&(60z\le \kappa),\\
 z(s)&=z_0+(\lambda-\kappa)s&&(60z\ge \kappa).
 \end{aligned}
 \label{eq:structure-queue-flows}
\end{equation}
The branches join at $z=\kappa/60$. Verified exponential and logarithmic
evaluation bounds their crossing times. If a forecast interval has an
unresolved common branch, its extremal offered arrivals give continuous
queue bounds. When those bounds establish one branch throughout the
interval and the forecast is $\ell_0+\ell_1s$, sharper endpoints are
\begin{equation}
\begin{aligned}
 z(h)&=e^{-60h}z_0+(\ell_0\pm\bar w_x)\varphi_1+\ell_1\varphi_2
       &&\text{(supply)},\\
 z(h)&=z_0+(\ell_0\pm\bar w_x-\kappa)h+\ell_1h^2/2
       &&\text{(gate)},\\
 \varphi_1&=(1-e^{-60h})/60,\qquad
 \varphi_2=(h-\varphi_1)/60.
\end{aligned}
\end{equation}
The minus and plus signs give lower and upper solutions. These sharpen
the endpoint bound while the established continuous enclosure remains
in force. Initial physical boxes must satisfy class and queue
nonnegativity. Subsequent endpoint sets may be intersected with these
independent bounds and then enclosed outward, including recentering error.

To retain correlations, on interval $l$ write
\begin{equation}
 \begin{aligned}
 L(x-p_l)&=\Gamma_l(s)\xi+\zeta_l(s),\qquad |\xi|\le\mathbf1,\\
 \Gamma_l(s)&=e^{A_l^0s}\Gamma_l(0),\qquad |\zeta_l(0)|\le\rho_l.
 \end{aligned}
 \label{eq:structure-affine-set}
\end{equation}
The same $\xi$ describes initial uncertainty across the plan. The fixed
matrix $A_l^0$ propagates its signed coefficients, and the remainder
$\zeta_l$ may depend on $\xi$. A fresh physical box
$x(t_0)\in\widehat x+[-d_x,d_x]$ initializes
$p(t_0)=\widehat x$, $\Gamma(0)=L\operatorname{diag}(d_x)$ and $\rho=0$
before numerical conversion. Conversion error is then added outward to $\rho$.
With mean-value matrix $B_l(s)$ and predictor defect
$\epsilon_l=Lf_b(p_l,t)-L\dot p_l$,
\begin{equation}
 \begin{aligned}
 \dot\zeta_l={}&B_l(s)\zeta_l+(B_l(s)-A_l^0)\Gamma_l(s)\xi\\
              &+\epsilon_l+(Le_7)w.
 \end{aligned}
 \label{eq:structure-remainder-identity}
\end{equation}
Every Jacobian bound must cover the full predictor-to-true-state segment.
For the queue coordinate, intersect the original trial enclosure with the
convex hull of predictor queue values and independent continuous true-queue
bounds. This retains the segment even when the predictor lies outside the
true queue range. A branch may be excluded only throughout this restricted
enclosure.

Omit $l$, and let $\mathcal J_\varpi$ contain the retained transformed branch
Jacobians. A matrix with nonnegative off-diagonal entries (a Metzler
matrix) and a forcing bound are
\begin{equation}
 \begin{aligned}
 \overline{M}_{ii}&=\max_\varpi\sup(\mathcal J_\varpi)_{ii},\\
 \overline{M}_{ij}&=\max_\varpi\sup|(\mathcal J_\varpi)_{ij}|\quad(i\ne j),\\
 (b_\zeta)_i&\ge\max_\varpi\sup_{J\in\mathcal J_\varpi,\,F\in\mathcal F}
                  \sum_j|[(J-A^0)F]_{ij}|\\
             &\hspace{1em}+\delta_i+|(Le_7)_i|\bar w_x,
 \end{aligned}
 \label{eq:structure-metzler}
\end{equation}
where $\mathcal F$ contains every $e^{A^0s}\Gamma(0)$ on $[0,h]$ and
$\delta_i\ge\sup_s|\epsilon_i(s)|$. The remainder comparison is
\begin{equation}
 |\zeta(s)|\le e^{\overline{M}s}\rho+\Psi(s)b_\zeta,
 \qquad \Psi(s)=\int_0^s e^{\overline{M}v}\,dv.
 \label{eq:structure-remainder-bound}
\end{equation}
The exponential is nonnegative and $\Psi$ is componentwise increasing.
Individual entries of the exponential may decrease, so the inherited
remainder must be bounded over the full interval. Use outward bounds
\begin{equation}
 \begin{aligned}
 f_i&\ge\sup_{s\in[0,h]}\sum_j|(e^{A^0s}\Gamma(0))_{ij}|,\\
 r_i^\zeta&\ge\sup_{s\in[0,h]}(e^{\overline{M}s}\rho)_i
                            +(\Psi(h)b_\zeta)_i.
 \end{aligned}
 \label{eq:structure-continuous-radii}
\end{equation}
Strict containment requires $f+r^\zeta<R$. For any row
$a_x^\top x\le b_x$, its continuous maximum is bounded by
\begin{equation}
 \sup_s a_x^\top p(s)+
 \sup_s\|a_x^\top L^{-1}e^{A^0s}\Gamma(0)\|_1+
 |a_x^\top L^{-1}|r^\zeta.
 \label{eq:structure-signed-output}
\end{equation}
The signed constraint coefficients are combined before taking absolute
values, preserving cancellations between correlated state errors. At a
scheduled endpoint use the endpoint generator and remainder. Independently
proved positivity, queue bounds, and exact constant-command rows may
establish their own constraints, so each row can use its tighter valid bound.

Let $\widehat\Gamma$ be the stored endpoint generator and let $d_\Gamma$
bound the row sums of $|e^{A^0h}\Gamma(0)-\widehat\Gamma|$. Recenter using
\eqref{eq:new-recentering-bound} and update
\begin{equation}
 \rho^+\ge d_\Gamma+e^{\overline{M}h}\rho+\Psi(h)b_\zeta+\delta^{\rm rec}.
 \label{eq:structure-endpoint-update}
\end{equation}
The same $\xi$ is retained. The total queue error combines generator and
remainder terms that may cancel, so a bound on their sum cannot be applied
separately to either term.

The next proposition verifies that these traffic-specific refinements
preserve the physical and service guarantees.

\begin{proposition}
\label{prop:structure-plan}
Consider a finite plan of nonnegative constant commands with nonnegative
initial destination-class inventories and queue, nonnegative class demands and
offered arrivals, and a nonnegative completion factor for all nonnegative
regional accumulations. Suppose that, on each plan interval, the positivity
and scalar queue bounds, signed uncertainty propagation, remainder comparison,
and strict self-inclusion test specified in this subsection all hold. If the
corresponding physical and scheduled service checks pass, then every admissible
trajectory satisfies the physical constraints throughout the plan horizon and
all included service requirements at their scheduled deadlines. Across
interval boundaries, recentering and generator-rounding errors are absorbed
into the propagated remainder while the signed initial-state dependence is
retained.
\end{proposition}

\begin{proof}
At every zero destination-class face, the class balances, including
\eqref{eq:structure-derived-classes}, have nonnegative inflow under the
stated sign assumptions. At $z=0$, admission is zero and the offered arrival
is nonnegative. For a local solution, the sum of negative class and queue
parts has upper derivative bounded by a local Lipschitz constant times
that sum. Its zero initial value remains zero by scalar comparison.
All classes and the queue are therefore nonnegative. Conservation gives
$0\le M(t)\le M(t_0)+\int_{t_0}^t(d^{\rm int}_1+d^{\rm int}_2+\Lambda(s))\,ds$.
Bounded arrivals prevent finite-time escape. These facts establish the
lower-state and class rows independently of upper storage constraints.

The scalar queue equation is Lipschitz in $z$ and increasing in the
arrival input. Scalar comparison orders solutions from ordered initial
queues and lower or upper arrival inputs, for every allowed measurable
residual. Equations~\eqref{eq:structure-queue-flows} and their verified
branch crossings thus give continuous queue bounds. The affine formulas
give valid endpoint refinements when one branch has already been proved.
Intersections preserve every reachable state. The predictor-to-state
segment remains inside the specified Jacobian region because its queue
coordinate lies in the convex hull of the two endpoint ranges.

Differentiating \eqref{eq:structure-affine-set}, using the mean-value
relation from Theorem~\ref{thm:new-plan-enclosure}, gives
\eqref{eq:structure-remainder-identity}. Its componentwise absolute
derivative satisfies $\frac{d}{ds}|\zeta|\le\overline{M}|\zeta|+b_\zeta$
almost everywhere. The bounds hold also for convex combinations of retained
branch Jacobians. Since $\overline{M}$ is Metzler, its exponential is
nonnegative. Variation of constants then proves
\eqref{eq:structure-remainder-bound}. Before a hypothetical exit from the
trial enclosure, \eqref{eq:structure-continuous-radii} bounds the complete
error by $f+r^\zeta<R$. Continuity rules out the exit. The independent
positivity and queue comparisons used here do not assume this containment.

For a row, substitution of the generator decomposition and
$|\xi|\le\mathbf1$ gives \eqref{eq:structure-signed-output}. The endpoint
form proves scheduled service checks. At a boundary, multiplying the generator conversion error by $\xi$
gives an error bounded by $d_\Gamma$. Adding this bound to the propagated
remainder and the predictor-recentering error yields
\eqref{eq:structure-endpoint-update}. This keeps the same initial
uncertainty variables and covers all conversion errors. Finite induction
proves containment and every physical and service row over the plan.
\end{proof}

Numerical evaluation of the matrix functions also needs verified error
bounds. For $\alpha\ge\|A\|_\infty$, $\sigma=\alpha h$, and series order $p$,
\begin{equation}
 \begin{aligned}
 \left\|e^{Ah}-\sum_{k=0}^p\frac{A^kh^k}{k!}\right\|_\infty
       &\le e^\sigma\frac{\sigma^{p+1}}{(p+1)!},\\
 \left\|\int_0^h e^{As}\,ds-
          \sum_{k=0}^p\frac{A^kh^{k+1}}{(k+1)!}\right\|_\infty
       &\le he^\sigma\frac{\sigma^{p+1}}{(p+2)!}.
 \end{aligned}
 \label{eq:structure-matrix-tails}
\end{equation}
The norm exponential series gives the first bound. Integration with
$e^{\alpha s}\le e^{\alpha h}$ gives the second. Signed row coefficients
are formed before interval-time evaluation. Outward rounding covers every
coefficient, tail, product, sum, and generator conversion used in the
preceding tests. Equivalent scalar evaluation shortcuts and regression
records are provided with the repository implementation.

\subsection{Plan replacement and repair}
\label{sec:new-execution}

At installation, the supervisor records the initial state set, model and
constraint versions, uncertainty bounds, service obligations, exact
commands, absolute timing, correlated reachable-set history, and expiry.
Rejection retains that record and its remaining schedule, including
already certified command changes. Restricting its verified trajectory
family to the remaining interval preserves the proof. A checkpoint box
$|L(x-p)|\le|\Gamma|\mathbf1+\rho$ contains reachable states but loses
correlations. Continuing the original certificate uses its stored history.
A new proof can instead use the correlated set or a containing box,
provided it verifies the entire chosen set.

A fresh trusted observation
$\widehat x_i=x(t_i)+\nu_i$, $|\nu_i|\le d_i$, supplies
$\mathcal I_i=\widehat x_i+[-d_i,d_i]$ at the actual installation time.
Its timestamp and error bounds are trusted inputs, and the whole set is
certified even if it is larger than the incumbent's reachable set. Without
a fresh observation, use an enclosure propagated under the incumbent's
schedule and uncertainty bounds. Intersections are valid when each set
contains the actual state. A changed forecast, error bound, receiving-supply
model, or task version requires a matching certificate. Every replacement
must cover all pending deadlines, and the incumbent's assumptions must
remain valid until replacement. Its finite guarantee ends at its original
expiry.

For a certified plan $P$, let $\tau_P$ be its installation time and
let $\mathcal X_P$ be the trajectory family covered by its stored record:
\begin{equation}
\mathcal X_P=\{x_P(\cdot;x_0,w):
 x_0\in\mathcal I_{\tau_P},\ w\in\mathcal W_P\}.
\label{eq:execution-covered-family}
\end{equation}
Its expiry is $E_P$, and $\mathcal O_P$ contains its certified service
triples $(T_r,g_r,\eta_r)$. Theorems and propositions on one-plan
certification establish the physical rows and these service rows for every
member of $\mathcal X_P$. Retention restricts this same family to the
remaining interval, keeping its correlated history. A new hard requirement
is admitted only when the plan retained or installed covers it. A soft
priority request alone does not change the hard contract.

The following theorem carries these one-plan guarantees through successive
replacement or retention decisions.

\begin{theorem}
\label{thm:new-plan-closure}
Consider $t_0<\cdots<t_N<t_{N+1}=T$, including every plan-admission
decision and admission of a new service requirement on this finite
interval. Let $\mathcal X_0$ be the initial certified family, whose initial
set contains the actual state. For a proposal $\widehat P_i$, write
$\widehat{\mathcal X}_i=\mathcal X_{\widehat P_i}$. At $t_i$, accept it only when
$\mathsf{Cert}_{t_i}(\widehat P_i\mid\mathcal S_{t_i})=1$, its initial
set contains $x(t_i)$, and it covers all pending admitted requirements.
The selected trajectory family is
\begin{equation}
\mathcal X_i=\begin{cases}
\widehat{\mathcal X}_i,&\text{if the replacement is accepted},\\
\mathcal X_{i-1},&\text{otherwise},
\end{cases}\qquad i=1,\ldots,N.
\label{eq:execution-family-update}
\end{equation}
Let $P_i,E_i,\mathcal O_i$ be the plan, expiry and service rows stored
with this family. Assume the executed model, uncertainty and commands
satisfy its certificate, and
\[
t_{i+1}\le E_i,\qquad \mathcal O_{t_i}\subseteq\mathcal O_i,
\qquad i=0,\ldots,N,
\]
where $\mathcal O_{t_i}$ contains the pending admitted requirements.
Then
\begin{equation}
\begin{aligned}
H_jx(t)+J_ju^{\rm exec}(t)&\le\beta_j^{\rm phys},
 &&t\in[t_0,T],\\
g_r^\top x(T_r)&\le\eta_r,
 &&\text{admitted rows with }t_0\le T_r\le T.
\end{aligned}
\label{eq:execution-physical-service}
\end{equation}
Rows and uncertainty bounds refer to the applicable active certificate.
The finite guarantee extends to its expiry when no later plan is installed.
\end{theorem}

\begin{proof}
The initial certificate and realized uncertainty place the executed
trajectory in $\mathcal X_0$ until the first decision. Suppose this holds
through $t_i$. Rejection preserves the same trajectory family, assumptions,
command schedule and original expiry. Its stored guarantee therefore
covers the next executed segment. Acceptance gives
$x(t_i)\in\mathcal I_{t_i}$ and a certificate for every allowed future
uncertainty realization, so the new segment belongs to the replacement
family. In either case,
\[
x|_{[t_i,t_{i+1}]}\in
\{\widetilde x|_{[t_i,t_{i+1}]}:\widetilde x\in\mathcal X_i\}.
\]
Induction over $i=0,\ldots,N$ proves the physical inequalities on each
segment. States are continuous at switches, and successor command rows are
checked at installation. Every admitted service row remains covered until
its deadline by $\mathcal O_{t_i}\subseteq\mathcal O_i$ and
$t_{i+1}\le E_i$. Applying its covering certificate at $T_r$ gives the
second inequality in \eqref{eq:execution-physical-service}.
\end{proof}

\subsubsection{Finite-grid plan repair}
\label{sec:predictive-repair}

Retention may forgo a useful requested priority. A predictive filter,
related to \cite{wabersich2021}, therefore tests interpolations between
the proposal and the certified continuation. Its common decision record is
\[
\mathcal S_t=(\mathcal I_t,\mathcal W_t,\mathcal O_t,P_t^B,t,\mathsf v_t),
\]
with current state information, uncertainty bounds, pending obligations,
backup plan, time, and model/constraint versions. All tested candidates
use this same record. Align constant command vectors
$P^0=(u_k^0)_{k=0}^{K-1}$ and $P^B=(u_k^B)_{k=0}^{K-1}$ on
$t=\tau_0<\cdots<\tau_K=E_B$, subdividing blocks if needed, and set
\begin{equation}
\begin{aligned}
P(\lambda_{\rm rep})&=(1-\lambda_{\rm rep})P^0+\lambda_{\rm rep} P^B,\\
\mathcal G_{\rm rep}&=\{0=\lambda_{{\rm rep},0}<\cdots<\lambda_{{\rm rep},N_{\rm rep}}=1\}.
\end{aligned}
\label{eq:repair-family}
\end{equation}
The grid has $N_{\rm rep}+1$ weights, ordered from the proposal at zero
to the backup at one.
Convex command constraints hold for interpolations of admissible endpoints.
State, queue, and deadline constraints require certification.
Let $\mathsf{Cert}_t(P\mid\mathcal S_t)=1$ mean that the exact plan is certified
for every specified initial state and uncertainty realization, continuously
for physical rows and at every pending service deadline over $[t,E_B]$.
Then
\begin{equation}
\lambda_{\rm rep}^*=\min\{\lambda_{\rm rep}\in\mathcal G_{\rm rep}:
                  \mathsf{Cert}_t(P(\lambda_{\rm rep})\mid\mathcal S_t)=1\}.
\label{eq:repair-selection}
\end{equation}
Test the proposal first, then every earlier grid point before selecting
a later one. Install a passing plan only when its installation conditions
match. Retain the incumbent on exhaustion or timeout. A prior verdict
applies only to the identical plan and decision record. The backup endpoint
must be certified for this same state set, including any enlarged
observation box.

The following theorem establishes termination and the smallest command
change among certified candidates on the tested grid.

\begin{theorem}
\label{thm:repair-grid}
Suppose $\mathsf{Cert}_t(P^B\mid\mathcal S_t)=1$, $1\in\mathcal G_{\rm rep}$, and each
grid evaluation terminates. Ordered exhaustive search returns a certified
plan after at most $N_{\rm rep}+1$ evaluations. For any norm on the stacked command
vector, the first certified grid point minimizes
$\operatorname{dist}(P,P^0)=\|P-P^0\|$ over the certified points of this grid.
\end{theorem}

\begin{proof}
Since $P(1)=P^B$, the certified grid set is nonempty. Finite ordered search
therefore terminates at its least element $\lambda_{\rm rep}^*$, and positive
homogeneity gives
\[
\operatorname{dist}(P(\lambda_{\rm rep}),P^0)=\lambda_{\rm rep}\|P^B-P^0\|.
\]
Every smaller grid point has been tested and rejected. The displayed
distance is nondecreasing in $\lambda_{\rm rep}$, including the case $P^B=P^0$,
which proves the stated minimum. The argument permits arbitrary certificate
patterns between successive grid points.
\end{proof}

This minimum is relative to the specified interpolation path, finite grid,
and sufficient checker. Minimization over all weights requires a further
analysis of the acceptance test. During computation the incumbent continues.

The next proposition verifies execution after repair and retention when
the decision deadline is missed.

\begin{proposition}
\label{prop:repair-execution}
Let $P^B$ be a valid incumbent on $[t,E_B]$, and consider an
installation time $t_a\in[t,E_B)$. A tested repair $P^R$ has expiry
$E_R>t_a$ and a completed-certificate time $t_{\rm ready}$, with
$t_{\rm ready}=\infty$ if no result is obtained. Assume the incumbent's
contract remains valid through preparation and installed records match the
actual time, model and coefficients. A repair installation requires
$x(t_a)\in\mathcal I_{t_a}$. All pending admitted service rows must be
covered by the selected record. Use
\begin{equation}
P^{\rm exec}=\begin{cases}
P^R,&t_{\rm ready}\le t_a\ \text{and}\quad
 \mathsf{Cert}_{t_a}(P^R\mid\mathcal S_{t_a})=1,\\
P^B,&\text{otherwise}.
\end{cases}
\label{eq:repair-installation-rule}
\end{equation}
The incumbent executes before $t_a$. Let $E_{\rm exec}$ be the selected
record's expiry. Then every physical row holds on $[t,E_{\rm exec}]$,
and every admitted service row in that record holds at its deadline.
Retention uses $E_{\rm exec}=E_B$ and the incumbent's original proof.
\end{proposition}

\begin{proof}
The incumbent certificate covers the prefix $[t,t_a]$. In the first
case of \eqref{eq:repair-installation-rule}, the certificate contains the
actual installation state and covers $[t_a,E_R]$. Concatenation gives
\eqref{eq:execution-physical-service} through $E_R$. In the second case,
the executed suffix remains a restriction of a trajectory in
$\mathcal X_{P^B}$, including its stored uncertainty dependence.
Its physical and service checks therefore remain valid through $E_B$.
The readiness condition prevents a late verdict from changing the command
already executed at $t_a$. A certificate prepared earlier is usable only
when it certifies this actual installation record.
\end{proof}

\subsection{Certification during computation}
\label{sec:planned-handover}

Choose a future activation $t_a\in(t,E_B)$. Let
$\mathcal R_B(t_a\mid\mathcal I_t)$ contain every state reachable from
$\mathcal I_t$ by the incumbent's original absolute schedule under its
allowed uncertainty. Certify a replacement from
$\mathcal I_{t_a}^+\supseteq\mathcal R_B(t_a\mid\mathcal I_t)$, with commands
starting at $t_a$ and remaining deadlines kept absolute. The incumbent
continues while this proof is prepared.

The following proposition covers an on-time handover and retention after
an unavailable or late replacement certificate.

\begin{proposition}
\label{prop:planned-handover}
Suppose the incumbent has a valid certificate on $[t,E_B]$ and the
actual state at $t$ belongs to both $\mathcal I_t$ and its certified
reachable set. A replacement is certified from $\mathcal I_{t_a}^+$ on
$[t_a,E_R]$, where $E_R>t_a$, with
$\mathcal R_B(t_a\mid\mathcal I_t)\subseteq\mathcal I_{t_a}^+$.
Every service obligation pending at $t_a$ has its deadline in $[t_a,E_R]$
and is included in the replacement certificate. Assume that the model,
uncertainty, obligations and incumbent schedule used in propagation
remain valid through $t_a$, and that the executed model and uncertainty
satisfy the replacement contract through $E_R$.
If the replacement certificate is ready by
$t_a$ and the authority installs exactly its future-dated schedule at $t_a$,
physical constraints and admitted service rows hold through $E_R$.
If the certificate is unavailable or late and the incumbent contract
remains valid through $E_B$, retaining its original schedule preserves
its guarantees through $E_B$.
\end{proposition}

\begin{proof}
The incumbent certificate covers the executed prefix $[t,t_a]$ and every
service row due on that prefix. Because propagation encloses every state reachable from the current
information set under the incumbent's schedule and uncertainty bounds,
the actual activation state lies in
$\mathcal R_B(t_a\mid\mathcal I_t)$ and hence in $\mathcal I_{t_a}^+$. The replacement-contract assumptions and exact
future-dated commands allow the new certificate to cover the replacement
trajectory and every remaining admitted service row on $[t_a,E_R]$.
Concatenating these two certified intervals gives the installed-plan
claim. If preparation is late or unsuccessful, the unchanged incumbent
continues the same previously certified trajectory under the incumbent
contract. Its reachable-set proof and expiry remain valid. This is the retention case of
Theorem~\ref{thm:new-plan-closure}.
\end{proof}

The activation enclosure may be intersected with a timestamped measurement
set that also contains the actual state. A checkpoint hull can initialize
a new proof. Retained execution still uses the original correlated proof.
Repair for the future slot uses the remaining common absolute partition
and checks its backup endpoint from that slot's information set. A task
or model change during preparation requires a fresh matching certificate.

The checker certifies the coefficients actually installed. The distance
identity in Theorem~\ref{thm:repair-grid} uses exact arithmetic. Stored
finite-precision grid plans must be ordered by their actual Euclidean
coefficient distances from the proposal, or checked to follow that order.
Minimum intervention then applies to the grid for the actual activation
record. Each installed plan has its own certificate for those stored
coefficients and that installation time.

\subsection{Networks with several entry queues}
\label{sec:network-extension}

Let $c\in\R^q$ contain destination-class inventories, $z\in\R^{n_q}$ exterior
queues, and $a\in\R^{n_q}$ admitted flows. For fixed allocation shares
$P_a\ge0$ with $\mathbf1_q^\top P_a=\mathbf1_{n_q}^\top$,
\begin{equation}
 \dot c=F(c,u,t)+P_a a,\qquad \dot z=\Lambda-a,
 \label{eq:network-class-queue}
\end{equation}
where $F$ contains internal demand, transfers, and trip completion. Regions
may have different production functions and routing patterns.

The following proposition extends conservation coordinates to several
queues while counting every internal and waiting vehicle exactly once.

\begin{proposition}
\label{prop:network-localization}
For fixed $P_a$, the invertible coordinate map
\begin{equation}
 L_a=\begin{bmatrix}I&P_a\\0&I\end{bmatrix},\qquad
 y=L_a\begin{bmatrix}c\\z\end{bmatrix}
   =\begin{bmatrix}c+P_a z\\z\end{bmatrix}
 \label{eq:network-transform}
\end{equation}
gives
\begin{equation}
 \dot y_c=F(y_c-P_a z,u,t)+P_a\Lambda,\qquad
 \dot z=\Lambda-a.
 \label{eq:network-transformed-flow}
\end{equation}
Thus explicit admission terms occupy only the $n_q$ queue coordinates. No
invertible linear transformation can confine these $n_q$ independent channels
to fewer than $n_q$ scalar coordinates. Moreover,
$\mathbf1_q^\top y_c=\mathbf1_q^\top c+\mathbf1_{n_q}^\top z$, so the combined
inventories count all internal and waiting vehicles exactly once.
\end{proposition}

\begin{proof}
Differentiate $y_c=c+P_a z$ with fixed $P_a$:
$\dot y_c=F(c,u,t)+P_a a+P_a(\Lambda-a)=F(c,u,t)+P_a\Lambda$.
Substituting $c=y_c-P_a z$ proves
\eqref{eq:network-transformed-flow}. The original admission incidence
$[P_a^\top,-I]^\top$ has rank $n_q$ because its lower block is $-I$.
Invertible coordinate changes preserve rank, so fewer than $n_q$ coordinates
cannot carry these independent channels. Finally,
$\mathbf1_q^\top y_c=\mathbf1_q^\top c+\mathbf1_q^\top P_a z
=\mathbf1_q^\top c+\mathbf1_{n_q}^\top z$ by the column-sum condition.
Thus the transformed inventories count every vehicle once.
\end{proof}

For independent queues and fixed block commands,
$a_i=\min\{k_i z_i,\kappa_i g_i\}$ with $k_i,\kappa_i\ge0$. With
$J_c=D_cF$ and diagonal admission derivative $D_a$,
\begin{equation}
 J_y=\begin{bmatrix}J_c&-J_cP_a\\0&-D_a\end{bmatrix},
 \qquad (D_a)_{ii}\in[0,k_i].
 \label{eq:network-jacobian}
\end{equation}
Queue branches affect only the diagonal entries. Regional coupling remains
in $J_c$ and $-J_cP_a$. Arrival residuals enter through
$[P_a^\top,I]^\top w$. The same Jacobian, generator, remainder, support,
and first-exit arguments apply in dimension $q+n_q$, with componentwise
residual bounds and the corresponding transformed disturbance incidence.
Separate queue-diagonal bounds are valid independently of simultaneous
branch combinations. The predictor defect still includes all branch
combinations not excluded by the local tests.

Forecast jumps at partition boundaries preserve physical states, so each
interval uses its affine forecast and inherits the endpoint set. This
extension assumes fixed allocation and independent admission. Changing
allocation adds derivative or mismatch terms, while shared receiving
capacity, state-dependent admission, or feedback gates require their actual
coupled derivatives. Those implementation extensions and the network
experiment specification are retained in the repository.

\subsection{Continuation beyond plan expiry}
\label{sec:terminal-backup-main}
\label{sec:terminal-backup}

A finite-plan certificate covers execution through its expiry $E$, which
is preserved after rejection. Coverage beyond $E$ requires a replacement
installed by $E$ with a later certified expiry or a certified terminal
backup activated at $E$. If neither is available, the finite-plan result
establishes feasibility only through $E$.

In this framework, indefinite continuation uses a nonempty certified set
$\mathcal K$ and an implementable feedback policy $\mu_B$. For every
activation state in $\mathcal K$ and every allowed uncertainty realization,
the backup trajectory exists for all future time, stays in $\mathcal K$,
and satisfies the physical state and command limits. Controller
initialization is part of the certificate.

The following proposition gives sufficient conditions for such a backup.

\begin{proposition}
\label{prop:backup-lyapunov}
Suppose the implemented backup has closed-loop dynamics
\[
\dot x=g_B(x,w_B),\qquad w_B(t)\in\mathcal W_B,
\]
where $\mathcal W_B$ is compact, $w_B(\cdot)$ is measurable, and $g_B$ is
continuous and locally Lipschitz in $x$, uniformly over
$w_B\in\mathcal W_B$, on an open domain $O$. Let
$V_B:O\rightarrow\mathbb R_{\ge0}$ be continuously differentiable and define
\[
\mathcal K=\{x\in O:V_B(x)\le\ell_B\},
\]
where $\mathcal K$ is nonempty and compact. If, on a neighborhood of
$\mathcal K$,
\[
\nabla V_B(x)^\top g_B(x,w_B)
\le -\lambda_B V_B(x)+\bar d_B
\]
for all $w_B\in\mathcal W_B$, with
$\lambda_B>0$ and
$0\le\bar d_B\le\lambda_B\ell_B$, and if every state and command generated
on $\mathcal K$ satisfies the authorized physical constraints, then
$\mathcal K$ is robustly invariant under $\mu_B$.
\end{proposition}

\begin{proof}
Along every admissible backup trajectory activated at time $t_a$,
\[
\frac{d}{dt}V_B(x(t))
\le -\lambda_B V_B(x(t))+\bar d_B
\]
almost everywhere while the trajectory remains in the stated neighborhood.
Scalar comparison gives
\[
V_B(x(t))
\le
e^{-\lambda_B(t-t_a)}V_B(x(t_a))
+
\frac{\bar d_B}{\lambda_B}
\bigl(1-e^{-\lambda_B(t-t_a)}\bigr).
\]
If $x(t_a)\in\mathcal K$, then
$V_B(x(t_a))\le\ell_B$ and
$\bar d_B/\lambda_B\le\ell_B$, so
$V_B(x(t))\le\ell_B$. The trajectory therefore cannot leave
$\mathcal K$. Because $\mathcal K$ is compact and contained in the domain of
the closed-loop dynamics, the solution extends for all future times. The
assumed state and command constraints then remain satisfied under the backup.
\end{proof}

A plan installed at $t_i$ with expiry $E_i>t_i$ is terminally certified
when it passes the physical and service checks on $[t_i,E_i]$, covers
every obligation pending at installation at a deadline in that interval,
and its entire terminal reachable set satisfies
\begin{equation}
\mathcal R_i(E_i)\subseteq\mathcal K.
\label{eq:backup-terminal-landing}
\end{equation}
Its certificate specifies an implementable transition to $\mu_B$ at
$E_i$, including required initialization. For
Proposition~\ref{prop:backup-lyapunov}, a sufficient test is
\begin{equation}
\sup_{x\in\mathcal R_i(E_i)}V_B(x)\le\ell_B.
\label{eq:backup-computable-landing}
\end{equation}
Write $\mathsf{Cert}^{\mathcal K}_t$ for the checker with these
terminal requirements. It returns one only after verifying finite-plan
checks, pending-service coverage, an implementable transition, and
\begin{equation}
\mathsf{Cert}_t(P\mid\mathcal S_t)=1,\qquad
\mathcal R_P(E_P)\subseteq\mathcal K.
\label{eq:terminal-checker}
\end{equation}
Here $\mathcal R_P$ is the certified reachable-state enclosure. Its
terminal inclusion covers every trajectory in $\mathcal X_P$.

The following theorem carries this terminal guarantee through successive
replacements and activates the backup when the active plan expires.

\begin{theorem}
\label{thm:backup-recursive-availability}
Let $(\mathcal K,\mu_B)$ be a certified terminal backup. Enumerate
actual plan installations as $t_0^{\rm ins}=t_0<t_1^{\rm ins}<\cdots$,
with plans $P_k$ and expiries $E_k$. Require
\[
\mathsf{Cert}^{\mathcal K}_{t_k^{\rm ins}}
 (P_k\mid\mathcal S_{t_k^{\rm ins}})=1,\qquad
x(t_k^{\rm ins})\in\mathcal I_{t_k^{\rm ins}},
\]
and preservation of all pending admitted service rows.
Assume $\#\{k:t_k^{\rm ins}\le T\}<\infty$ for every finite $T$, and
that execution and uncertainty satisfy the active certificate. Define
$k(t)=\max\{k:t_k^{\rm ins}\le t\}$ and execute
\begin{equation}
u^{\rm exec}(t)=\begin{cases}
u_{P_{k(t)}}(t,x(t)),&t<E_{k(t)},\\
\mu_B(x(t)),&t\ge E_{k(t)}.
\end{cases}
\label{eq:terminal-execution-policy}
\end{equation}
Use the certified controller initialization at expiry, and admit any new
service request during backup operation only through a covering terminally
certified plan. Then the trajectory exists for every $t\ge t_0$, and
\begin{equation}
\begin{aligned}
H_jx(t)+J_ju^{\rm exec}(t)&\le\beta_j^{\rm phys},
 &&t\ge t_0,\\
g_r^\top x(T_r)&\le\eta_r,
 &&\text{admitted rows with }T_r\ge t_0.
\end{aligned}
\label{eq:terminal-global-guarantee}
\end{equation}
Rejection preserves the active plan and its terminal transition.
\end{theorem}

\begin{proof}
Until a plan's expiry or replacement, Theorem~\ref{thm:new-plan-closure}
gives its finite physical and service guarantees. If
$E_k<t_{k+1}^{\rm ins}$, no replacement occurs before expiry, so
\[
x(E_k)\in\mathcal R_{P_k}(E_k)\subseteq\mathcal K
\quad\Longrightarrow\quad
x(t)\in\mathcal K\quad(E_k\le t<t_{k+1}^{\rm ins}).
\]
The implication follows from the backup's robust invariance and certified
initialization. Its state and command limits hold throughout this interval.
If no next installation exists, take $t_{k+1}^{\rm ins}=\infty$.
Every accepted plan starts a new certified finite segment. A rejected
request preserves the current segment and terminal connection. On each
bounded time interval there are finitely many installations, hence finitely
many such concatenations. Each finite plan exists through its certified
horizon, and each backup segment exists indefinitely. These compatible
solutions therefore define a trajectory for all future time. Every
admitted service deadline is retained in a covering finite certificate.
At expiry no such deadline remains uncovered. Applying the covering
certificate at that deadline proves the second part of
\eqref{eq:terminal-global-guarantee}.
\end{proof}

A service request made during backup operation obtains execution authority
only after a new certificate covers that request and a terminal
continuation. A nonempty invariant set also requires arrivals compatible
with sustained discharge: finite exterior storage cannot accommodate
persistent allowed inflow above every attainable discharge rate. Any
diversion, upstream holding, or demand management used to meet this
condition must appear in the certified model. The numerical studies test
finite-plan validity. Indefinite operation has the additional terminal
conditions just proved.

The next corollary combines terminal continuation with ordered repair.

\begin{corollary}
\label{cor:repair-recursive}
Use the terminal checker in \eqref{eq:terminal-checker} for repair
at the actual installation record. Suppose
$\mathsf{Cert}^{\mathcal K}_t(P^B\mid\mathcal S_t)=1$ and
$1\in\mathcal G_{\rm rep}$. Then the terminally certified grid set
\begin{equation}
\mathcal G_{\rm rep}^{\mathcal K}(t)=
\{\lambda_{\rm rep}\in\mathcal G_{\rm rep}:
\mathsf{Cert}^{\mathcal K}_t(P(\lambda_{\rm rep})\mid\mathcal S_t)=1\}
\label{eq:repair-terminal-grid}
\end{equation}
is nonempty. Under the initialization, uncertainty, installation and
local-finiteness assumptions of
Theorem~\ref{thm:backup-recursive-availability}, installing only members
of this set and retaining the incumbent on timeout preserves
\eqref{eq:terminal-global-guarantee}.
\end{corollary}

\begin{proof}
The interpolation endpoint is $P(1)=P^B$, so
$1\in\mathcal G_{\rm rep}^{\mathcal K}(t)$. Every selected member
has the finite service coverage and terminal inclusion
\[
\mathcal R_{P(\lambda_{\rm rep})}(E_B)\subseteq\mathcal K.
\]
An exhaustive terminating search selects its least certified weight by
Theorem~\ref{thm:repair-grid}. If computation is stopped, retention
preserves the incumbent's terminal certificate and original expiry.
Both cases satisfy the installation rule of
Theorem~\ref{thm:backup-recursive-availability}, which gives the claimed
continuation and service guarantees.
\end{proof}

\subsection{Cost of finite-plan repair}
\label{sec:repair-cost-bound}

Consider two fixed command plans with the same initial state and
disturbance on $[t,t+H]$. Suppose both trajectories remain in one compact
domain on which, uniformly over allowed uncertainty,
\[
\|f(x,u,w)-f(y,v,w)\|
\le L_x\|x-y\|+L_u\|u-v\|.
\]
For a common running cost $\ell$ and terminal cost $g$, assume
$|\ell(x,u)-\ell(y,v)|\le L_{\ell x}\|x-y\|+
L_{\ell u}\|u-v\|$ and $|g(x)-g(y)|\le L_g\|x-y\|$.
Set $J(P)=\int_t^{t+H}\ell(x_P,u_P)\,ds+g(x_P(t+H))$.

The following proposition bounds the cost change caused by interpolation
toward the certified backup.

\begin{proposition}
\label{prop:repair-cost}
Under these assumptions, for the Euclidean norm on the stacked fixed-block
command vector,
\begin{equation}
|J(P(\lambda_{\rm rep}))-J(P^0)|
\le L_J\lambda_{\rm rep}\|P^B-P^0\|_2,
\label{eq:repair-cost-bound}
\end{equation}
where
\begin{align*}
L_J&=L_{\ell x}L_u B_H+L_{\ell u}H+L_gL_u A_H,\\
A_H&=\begin{cases}(e^{L_xH}-1)/L_x,&L_x>0,\\H,&L_x=0,\end{cases}\\
B_H&=\begin{cases}(e^{L_xH}-1-L_xH)/L_x^2,&L_x>0,\\H^2/2,&L_x=0.\end{cases}
\end{align*}
\end{proposition}

\begin{proof}
The maximum instantaneous command difference is at most
$\delta=\lambda_{\rm rep}\|P^B-P^0\|_2$. Gronwall's inequality gives
$\|x_{P(\lambda_{\rm rep})}(t+s)-x_{P^0}(t+s)\|\le L_u\delta A_s$.
Integrating the running-cost bound and adding the terminal bound yields
\eqref{eq:repair-cost-bound}, since $\int_0^H A_s\,ds=B_H$.
\end{proof}

The domain must contain both trajectories, including any rejected-reference
portion outside the physical constraint set. This is an additional
assumption to certification of the repaired plan. The locally Lipschitz
polynomial class flows and admission minimum give finite model constants
on such a common compact domain. Linear TTS and fixed-priority costs give
finite cost constants. If the same constants cover every paired initial
state and residual, the bound is uniform and also bounds worst-case costs
by $|\sup J_1-\sup J_0|\le\sup|J_1-J_0|$. A numerical $L_J$ requires
verified bounds on those constants for that domain. The ensemble separately
reports realized costs for its specified policies and priority weights.

\section{Numerical Evaluation}
\label{sec:new-evidence}

The studies test whether certification preserves useful plans, whether
traffic structure improves the bounds, and whether repair and retention
support successive requests, computation delays and agent errors.
The primary model, objectives and constraints are defined in
Section~\ref{sec:benchmark-formulation} and Section~\ref{app:new-model}.
The seven-state experiments use 6000~veh/h base exterior arrivals,
10,000-vehicle regional limits and 2000-vehicle boundary storage. Arrival
events increase the known mean by 20--50\% for 30--60 minutes, including
five-minute rise and recovery ramps. Certificates cover the complete
specified initial set and every measurable bounded arrival residual.
Matched replays evaluate selected initial states and residual profiles.
\appref{app:experimental-settings} gives the essential study settings.
Full reproduction instructions and extended comparisons are in the
\reproductiondetails{GitHub reproduction guide}.

\subsection{Unchanged plans and tighter traffic bounds}
\label{sec:new-incident-protocol}
\label{sec:controller-benchmarks}
\label{sec:structure-results}
\label{sec:new-mpc-results}
\label{sec:structure-online}
\label{sec:network-transfer-results}

Task Q compares economic MPC with the queue/exit extension of the
almost-smooth control-Lyapunov-function (CLF) controller \cite{geroliminis2013mpc,zhong2018robust}.
Its twelve contracts combine 20, 30, 40 and 50\% arrival increases with
30-, 45- and 60-minute durations, one-vehicle state half-widths and a
300~veh/h residual bound. All twelve MPC plans pass certification and
execute unchanged. Across 960 matched feasible comparisons
($12$ contracts, $16$ initial-state corners and $5$ residual profiles),
the median paired percentage reduction in TTS is 24\% (range 23--25\%)
relative to the specified CLF extension. Certification therefore preserves the benefit of these plans
while adding the stated uncertainty coverage and expiry. Additional
regulation and eight-event compatibility results are provided in
\reproductiondetails{the controller comparisons}.

The structure comparison holds the plan, forecast and 15-second proof
partition fixed. For one feedback plan, destination-class coordinates
give a Region-1 deadline bound of 2220 vehicles. Incorporating waiting
inventory lowers it to 2169, while the raw-accumulation enclosure remains
unresolved. Table~\ref{tab:structure-frontier} compares three enclosures
for a fixed two-block MPC plan. Regional/class half-width is $d$ vehicles
and queue half-width is $\min(d,10)$. Signed correlations extend the
largest tested certifiable regional width from 10 to 50 vehicles at a
60~veh/h residual. Positivity/queue bounds perform better for the
one-vehicle box with the larger residual, so the refinements complement
each other. Acceptance requires every physical and service constraint
to pass its certified bound check.

\begin{table}[t]
\centering\small
\caption{Certified Region-1 inventory bounds at minute 60 for a fixed
model predictive control (MPC) plan over minutes 0--60. Entries are rounded
upward in vehicles. The service limit is 2850. Here $d$ is the sensing
half-width of the first six state coordinates in \eqref{eq:new-plant}
(vehicles), the queue half-width is
$\min(d,10)$ vehicles, and $\bar w_{\rm veh}$ bounds the exterior-arrival
forecast residual in vehicles per hour (veh/h). Base interval uses the
componentwise enclosure in conservation coordinates. Positivity/queue adds
independently proved nonnegativity and scalar queue bounds. Signed also retains
shared initial-state dependence through signed generators and a bounded
remainder. A dash means the complete-plan certificate is unresolved.}
\label{tab:structure-frontier}
\setlength{\tabcolsep}{5pt}
\begin{tabular}{@{}rrrrr@{}}
\toprule
$d$ (veh) & $\bar w_{\rm veh}$ (veh/h) & Base interval & Positivity/queue & Signed\\
\midrule
1 & 60 & 1780 & 1729 & 1722\\
10 & 60 & 1845 & 1777 & 1725\\
50 & 60 & --- & --- & 1853\\
100 & 60 & --- & --- & ---\\
1 & 300 & --- & 2009 & 2256\\
25 & 300 & --- & --- & 2504\\
\bottomrule
\end{tabular}
\end{table}

For a 1700-vehicle service limit under one-vehicle sensing half-widths
and a 60~veh/h residual bound, the base interval and positivity/queue
methods first certify a nominal target of 1600 vehicles on the six-target
menu. Signed generators also certify 1625, with a deadline bound below
1689. The required margin falls from 100 to 75 vehicles and nominal TTS
falls by 0.90\% (41 veh-h) over minutes 0--60. The margin is
the difference between the authorized service limit and the target imposed
on the nominal forecast trajectory. The extended online comparison, target menu
and adverse realization are in \reproductiondetails{the service-margin study}.

The network construction also certifies fixed plans on connected two- and
three-region rings with six and twelve states. At regional half-width
$d=25$ vehicles, combined inventories certify the two-region plan with a
300~veh/h residual per queue and the three-region plan with a
60~veh/h residual. Their largest regional deadline bounds are 1552
and 2060 vehicles against a 2700-vehicle limit. Class coordinates leave
both sufficient proofs unresolved. Successful twelve-state proofs take
5.20--5.28 seconds for 120 intervals. The essential ring setup is in
\appref{app:network-transfer}. The full sixteen-attempt sweep and runtime
records are in \reproductiondetails{the connected-network study}.

\subsection{Deadline coverage and checked repair}
\label{sec:repair-results}

Task D fixes the 30\% one-hour event and minute-20 traffic state. The
72 development contracts cross state half-widths $d\in\{1,10,25,50\}$
vehicles, residual bounds $\bar w_{\rm veh}\in\{60,150,300\}$ veh/h,
deadline limits $\eta\in\{2750,2850,2950\}$ vehicles and Region-2 weights
$p\in\{2,4\}$. Eight held-out contracts use $d\in\{5,20\}$,
$\bar w_{\rm veh}\in\{100,225\}$, $\eta\in\{2800,2900\}$ and $p=3$.
Each constant plan has 640 replays: all 128 corners crossed with five
residual profiles. Full-scenario MPC constrains these scenarios and the
center forecast with one-vehicle prediction clearance.

\begin{table}[t]
\centering\small
\caption{Deadline-constrained Task D: 72 development and eight held-out
contracts over minutes 20--60. Entries count passing contracts out of 80.
Sample passes satisfy physical and service limits in all 640 replays
(128 initial-state corners times five arrival-residual profiles).
Family passes certify continuous physical limits and the minute-60 service
bound for every allowed initial state and residual. MPC means model predictive control.
Nominal MPC plans from the state-set center with zero residual.
Five-scenario MPC adds the all-low/all-high states with both constant
residual extremes. Full-scenario MPC uses all 640 combinations plus the
center forecast. All minimize center-forecast priority cost. Checked-repair
entries count all certified outputs: unchanged proposals and checked
interpolations toward a backup. Reserve refresh constructs a new certified
backup for nominal-plan repair when the old continuation fails its check.}
\label{tab:benchmark-deadlines}
\begin{tabular}{@{}lrr@{}}
\toprule
Planning/execution option & Sample passes & Family passes\\
\midrule
Nominal MPC & 36/80 & 36/80\\
Five-scenario MPC & 40/80 & 40/80\\
Full-scenario MPC & 80/80 & 44/80\\
Nominal MPC with checked repair & 74/80 & 74/80\\
Full-scenario MPC with checked repair & 74/80 & 74/80\\
Repair with checked reserve refresh & 80/80 & 80/80\\
\bottomrule
\end{tabular}
\end{table}

Table~\ref{tab:benchmark-deadlines} separates passing the replay grid from
certifying the full uncertainty family. Full-scenario MPC passes every
sampled contract but certifies 44/80. Fresh checking certifies the existing
continuation in 68/80. Ordered repair produces certified plans in 74/80.
New reserves certify in all twelve contracts where the old continuation
fails its check and enable repair in the remaining six. This tests the
same-contract backup premise of Theorem~\ref{thm:repair-grid}.
The physical control-barrier-function quadratic program (CBF-QP)
separately satisfies physical constraints in 80/80
adverse runs and the deadline in 44/80, with one adverse realization per
contract. Its scheduled-service performance is evaluated separately from the
instantaneous physical barriers.

Table~\ref{tab:benchmark-conflict} uses the common
$(d,\bar w_{\rm veh},\eta,p)=(1,300,2850,4)$ contract at the all-high
corner, where every initial state component is at its upper sensing bound,
with the arrival residual held at its positive upper bound. Nominal MPC violates storage and service.
The physical CBF protects storage but misses the deadline. Retention and
checked repair satisfy both. Full-scenario MPC passes this realization,
while its sufficient family certificate remains unresolved.

\begin{table}[t]
\centering\footnotesize
\caption{Matched adverse realization, minutes 20--60. TTS is total
time spent in the regions and exterior queue. $J_D$ weights Region~2 by
$p=4$. Both use veh-h. Costs and counts are rounded to whole
units. Peak queue is the maximum exterior count $100z$. $100n_1(E)$ is the
Region-1 count at deadline $E=60$ minutes. Their limits are 2000 and 2850
vehicles. P/S lists physical/service outcomes: P=pass, F=fail. Cert.
reports full-family certification: Yes=certified, No=unresolved,
dash=not assessed by that checker.
MPC means model predictive control, and CBF means control barrier function.
Retained continuation follows the previous certified plan. Nominal-plan
and Full-scenario repair interpolate their MPC proposals toward it.}
\label{tab:benchmark-conflict}
\setlength{\tabcolsep}{3pt}
\begin{tabular}{@{}lrrrrcc@{}}
\toprule
Policy & TTS & $J_D$ & Peak queue & $100n_1(E)$ & P/S & Cert.\\
\midrule
Nominal MPC & 3300 & 6328 & 2180 & 2852 & F/F & No\\
Physical CBF & 3295 & 6376 & 1844 & 2917 & P/F & ---\\
Retained continuation & 2695 & 6814 & 194 & 1727 & P/P & Yes\\
Nominal-plan repair & 3097 & 6427 & 1683 & 2310 & P/P & Yes\\
Full-scenario MPC & 3267 & 6357 & 1999 & 2848 & P/P & No\\
Full-scenario repair & 3075 & 6453 & 1547 & 2314 & P/P & Yes\\
\bottomrule
\end{tabular}
\end{table}

For nominal, retained and repaired plans $P^0,P^B,P^R$ on a common
realization $\omega$, percentage reduction in priority cost $\Delta J_\omega$, percentage change in queue-inclusive delay $\Delta D_\omega$ and
recovered requested benefit $B_\omega$ are
\begin{equation}
\begin{aligned}
\Delta J_\omega&=100\frac{J_D(P^B;\omega)-J_D(P^R;\omega)}{J_D(P^B;\omega)},\\
\Delta D_\omega&=100\frac{J_Q(P^R;\omega)-J_Q(P^B;\omega)}{J_Q(P^B;\omega)},\\
B_\omega&=\frac{J_D(P^B;\omega)-J_D(P^R;\omega)}{J_D(P^B;\omega)-J_D(P^0;\omega)}.
\end{aligned}
\label{eq:benchmark-repair-metrics}
\end{equation}
Positive $\Delta J_\omega$ means a reduction in priority cost relative to retention.
Positive $\Delta D_\omega$ means an increase in queue-inclusive delay. Both changes are percentages. The recovered
benefit is reported as $100B_\omega$ percent.
Across 24,320 matched realizations from 38 original-continuation repair
contracts, every benefit denominator is positive. Median priority
improvement is 5.7\%, delay increase 16.1\%, and recovered benefit 91.7\%.
Repair recovers much of the requested Region-2 preference at a measured
delay cost. In the displayed conflict, retention has the lowest TTS.
Full per-realization costs are in \reproductiondetails{the repair study}.

\subsection{Successive requests from the inherited traffic state}
\label{sec:sequential-handover-results}

A certified Task-D MPC repair with Region-2 weight $p=4$ starts at minute
20. Stored planner outputs request the $p=2$ full-scenario suffix at minute
25 and the $p=4$ nominal suffix at minute 30. This deterministic sequence
isolates execution admission from language interpretation. Physical limits,
uncertainty and the minute-60 deadline remain fixed. Each replacement is
checked from a box enclosing the active plan's reachable state. Rejection
retains its schedule, correlated proof and original expiry.

\begin{table}[t]
\centering\small
\caption{Successive changes in deadline-constrained Task D. Case entries
are stored contract identifiers. The 25/30 column gives the decisions at
minutes 25 and 30: A=accept the replacement, R=retain the current plan.
Unchecked feasible and Certified feasible count replays satisfying both
physical limits and the minute-60 service requirement out of 640 matched
initial-state/arrival realizations. TTS means total time spent, including
the exterior queue, over minutes 20--60, in veh-h.
Median TTS F/U/C gives three medians in order: F=retain the fixed initial
certified plan, including its scheduled command changes; U=execute both
requested replacements without checking; C=execute only certified
replacements and retain the current plan after rejection. Medians use all
640 replays, including infeasible unchecked runs. D015, D027 and D071 were
preselected. D025 is the first storage-violation case found in the complete
matched-pair scan.}
\label{tab:sequential-handover}
\setlength{\tabcolsep}{4pt}
\begin{tabular}{@{}lcccc@{}}
\toprule
Case & 25/30 & Unchecked feasible & Certified feasible & Median TTS F/U/C\\
\midrule
D015 & A/A & 640/640 & 640/640 & 3196/3081/3081\\
D027 & A/R & 640/640 & 640/640 & 3037/3081/2717\\
D071 & R/R & 640/640 & 640/640 & 3169/3094/3169\\
D025 & A/R & 512/640 & 640/640 & 2963/3070/2719\\
\bottomrule
\end{tabular}
\end{table}

For D015, D027, D071 and D025, the contract triples
$(d,\bar w_{\rm veh},\eta)$ are $(1,60,2850)$, $(1,300,2850)$,
$(50,60,2950)$ and $(1,300,2750)$, respectively, in vehicles, veh/h
and vehicles. Table~\ref{tab:sequential-handover} shows acceptance of both changes,
acceptance followed by retention, and retention of the initial plan.
D027's second request fails the sufficient service enclosure although
the sampled unchecked schedule passes. D071's propagated boxes include
inadmissible states. The full development scan contains 36 matched weight
pairs, 30 with the certified source plans needed for this sequence.
Unchecked switching overfills storage in three of these 30 cases, all
tested by the sequential checker.

In the first overflow case, D025, the minute-25 replacement passes and
the minute-30 request fails the queue check. Unchecked execution reaches
2005.59 vehicles against the 2000-vehicle limit, confirmed by independent
adaptive integration. Retaining the accepted suffix gives 640/640
physically and service-feasible replays, with median TTS 8.2\% below
retaining the initial plan. Its median changing-priority cost is 5.2\% above
unchecked switching. These tests assume installation at the specified
times. The following study includes computation delay. The other overflow
cases and complete records are in \reproductiondetails{successive requests}.

\subsection{Computation delays and agent-context errors}
\label{sec:delayed-handover-results}
\label{sec:task-state-live}
\label{sec:queue-live-pilot}
\label{sec:primary-memory-pressure}

The handover study tests Proposition~\ref{prop:planned-handover} on seventeen
contracts with one-, two- and five-minute preparation slots: 51 cases.
The accepted plan continues while its reachable set is propagated to the
slot. Replacements are checked from that future box through the unchanged
deadline. Late or unsuccessful results leave the prior schedule and
certificate in force through its expiry. There are 30 timely replacements
and 21 retentions. All 32,640 replays satisfy physical and service limits.
Three one-minute results are late. A longer slot gives more computation
time but changes the reachable box and remaining service horizon. A new
box-based proof can fail after class correlations are discarded.

\begin{table}[t]
\centering\small
\caption{Delayed activation and agent decisions. Top: Handover slot is
preparation time after minute 20. Replacements and Retentions count timely
certified installations and continued execution of the prior plan,
respectively (17 contracts per slot). Each row has
$17\times640=10{,}880$ replays through minute 60.
Middle: the interpretation pilot has 24 calls per interface/memory pair.
History only supplies observations and interaction memory. Task record
also supplies the active priority and its expiry. Clean memory is
unaltered. Poisoned memory contains misleading instructions. Objective
correct means agreement with the active priority. Tool correct means
agreement with the saved certified-tool selection rule. Bottom: each
interface has four traces and $4\times640=2560$ replays. Objective correct
assesses initial requests. Repairs executed counts traces installing a
certified grid repair. Feasible trajectories satisfy physical and service
constraints.}
\label{tab:delayed-handover}
\label{tab:benchmark-api-pilot}
\label{tab:benchmark-api-execution}
\setlength{\tabcolsep}{4pt}
\begin{tabular}{@{}llrr@{}}
\toprule
Handover slot & Replacements & Retentions & Feasible trajectories\\
\midrule
1 minute & 8/17 & 9/17 & 10,880/10,880\\
2 minutes & 11/17 & 6/17 & 10,880/10,880\\
5 minutes & 11/17 & 6/17 & 10,880/10,880\\
\midrule
Pilot interface & Memory & Objective correct & Tool correct\\
\midrule
History only & Clean & 19/24 & 22/24\\
History only & Poisoned & 15/24 & 24/24\\
Task record & Clean & 24/24 & 24/24\\
Task record & Poisoned & 24/24 & 24/24\\
\midrule
Actual planning & Objective correct & Repairs executed & Feasible trajectories\\
\midrule
History only & 1/4 & 1/4 & 2560/2560\\
Task record & 4/4 & 4/4 & 2560/2560\\
\bottomrule
\end{tabular}
\end{table}

The agent study uses \texttt{gpt-5.4-mini-2026-03-17}, reasoning disabled,
and strict function schemas. The 96-call pilot crosses three availability
situations, four operator histories, clean or corrupted memory, two
interfaces and two repetitions. It routes precomputed plans. Eight later
traces instead use the agent's objective in a fresh MPC solve and request
repair or retention after rejection. These pair four histories across
both interfaces under corrupted memory, with thirteen live calls and
640 replays per executed schedule.

The task record supplies the authoritative priority, its expiry and
interruption status. Corrupted memory requests a changed priority,
increased storage limit or unchecked execution. Table~\ref{tab:delayed-handover}
reports objective/tool correctness separately from traffic feasibility.
The task record improves objective selection in these small paired
examples. In three history-only traces, an accepted MPC plan satisfies the
physical and service constraints but uses the wrong regional objective.
Only one trace executes a checked repair. Certification enforces the
original physical and service limits
on every accepted plan. The pilot and actual-planning traces are separate
stages, evaluated through the accepted plans' finite validity intervals. Essential histories,
hard limits and scoring are summarized in \appref{sec:supp-execution}.
Literal prompts, replies and replay instructions are in
\reproductiondetails{the agent experiments}.

\FloatBarrier

\section{Operational Insights and Implementation}
\label{sec:new-discussion}
\label{sec:information-model-scope}

The experiments show how certification can preserve a requested improvement,
repair a proposal at a measured traffic cost, or retain an earlier plan when
a replacement fails. Tighter bounds increase the plans that can be accepted
under the same information and service requirements. Wider uncertainty may
require a freshly certified reserve.

The studies use a literature MFD, synthetic internal demands and known
event means. Certificates cover the stated sensing and exterior-arrival
uncertainty while production, routing, internal demand, release rate and
gate capacity remain fixed. Deployment requires an estimator that supplies
a joint set containing the actual destination-class inventories and queues,
including measurement bias and timing errors, together with arrival-error
bounds containing the actual residuals. Regional totals alone leave
destination mixtures undetermined. Additional model uncertainty requires
validated bounds on its effect throughout each proof interval. Fixed
network allocation assumptions are stated in
Section~\ref{sec:network-extension}. Implementation details describe
extensions for changing allocation.

Trajectory replays measure performance on selected realizations. The
certificate establishes feasibility for the full specified family. An
unresolved certificate therefore leaves feasibility undecided unless a
violation is independently exhibited. The queue-capped structure sweep and
all-state deadline sweep use different information sets. The connected-ring
study tests the network construction under its own routing assumptions,
and the agent study is a small demonstration with one recorded model
version. Runtime depends on hardware and load.

The experiments cover finite plans through their recorded expiries.
Indefinite continuation through the terminal backup additionally requires
the invariant policy, terminal-set inclusion and obligation coverage in
Section~\ref{sec:terminal-backup}. Sustained arrivals must be compatible
with available storage and discharge. Implementations also need a certified
initial plan, certificates matching the current model and task versions,
and enough time for propagation, planning, checking and actuation before
the active plan expires. The paper includes the complete mathematical
argument and essential experimental settings. The
\reproductiondetails{repository guide} supplies commands, solver settings,
saved plans and proof records.

\section{Conclusions}
\label{sec:new-conclusion}

This paper develops a framework for checking agent-proposed urban
perimeter-control plans before execution. It verifies physical constraints
and service deadlines within the stated traffic model and uncertainty
bounds, using vehicle conservation and the links between boundary queues and
regional traffic to tighten bounds on future states. Plans that pass
certification can be executed unchanged, and checked repair searches for
certified alternatives. Rejected or late replacements leave the previously
accepted plan in use through its certified validity interval. When the
required terminal conditions are also certified, feasible operation can
continue beyond plan expiry. Numerical studies show reduced conservatism and
recovery of feasible plans through repair. Further experiments examine
successive plan changes, connected networks, computation delays, and errors
in the information given to the agent, with certified execution maintained
in the tested cases.

Future work will focus on estimating vehicle counts by region and
destination from traffic measurements and calibrating the associated
uncertainty bounds. Another direction is to establish uncertainty bounds
that capture differences between the traffic model and observed network
behavior. We will also study when new plans should take effect, accounting
for sensing, computation, and actuation time, traffic uncertainty, and
remaining service deadlines. These developments will support evaluation of
the framework in operational traffic networks.

\FloatBarrier
\section*{CRediT authorship contribution statement}
\textbf{Ali Eslami:} Conceptualization, Methodology, Software, Validation,
Formal analysis, Investigation, Data curation, Visualization,
Writing and editing.
\textbf{Jiangbo Yu:} Supervision, Review and editing.

\section*{Declaration of competing interests}
The authors declare that they have no known competing financial interests or
personal relationships that could have appeared to influence the work
reported in this paper.

\section*{Data and code availability}
Reproducibility materials include the traffic models, planning and
certification implementations, experiment protocols, stored agent requests
and responses, proof records, and trajectory-replay outputs described in the
experimental appendices. The materials are maintained at
\url{https://github.com/AlienEslami/agentic-perimeter-runtime-certification}.


\appendix
% Use appendix letters for equations while retaining full appendix section titles.
\renewcommand{\theequation}{\Alph{section}.\arabic{equation}}
\section{Essential Experimental Settings}
\label{app:experimental-settings}
\label{app:reusable-benchmark}
\label{sec:supp-queue-settings}
\label{app:benchmark-details}

\paragraph{Primary state and plans}
Tasks Q and D start at minute 20 with the state center, in the order of
\eqref{eq:new-plant} and units of 100 vehicles,
\[
\begin{aligned}
\widehat x_{20}=(&21.020544601495207,23.156747836310867,\\
&14.209387548878382,14.877124601919054,\\
&2.4028346640030684,3.0184193919155455,\\
&1.2410781134202304).
\end{aligned}
\]
The original continuation uses $(1,1,0.667015677,1)$ in command order
$(u_{12},u_{21},u_{02},u_{20})$, in five- and
35-minute blocks ending at minute 60. Changed requirements trigger fresh
checking. If necessary, reserve construction tries network TTS followed
by moderate Region-2 priority. A repair endpoint must pass the same
contract. Repair tests $\mathcal G_{\rm rep}=\{j/20:0\le j\le20\}$ in order,
from the proposal at weight zero to the checked backup at weight one. Plans are
compiled and checked against exact actuator constraints. State/service
proof intervals are 15 seconds.

\paragraph{Replay and comparisons}
The five residual profiles are zero, constant positive, constant negative,
five-minute alternating and ten-minute sinusoidal, within each contract's
bound. Constant plans use all 128 initial-box corners. The CLF extension
uses sixteen corners, every eighth entry of the lexicographic
$\{-1,+1\}^7$ list, with the same profiles. Performance comparisons match
corner and residual labels. Replays use fourth-order Runge--Kutta
(RK4) integration with ten-second steps, supplemented by independent
adaptive eighth-order Runge--Kutta (DOP853) integration and conservation
checks. The CLF extension uses a regional-error correction with
linear-quadratic regulator (LQR) admission/exit feedback, reference
totals $(2700,3200)$ vehicles, state weights $(4,1,0.1,0.1,0.1,0.1,2)$
and input weight $3000I_4$, where $I_4$ is the four-by-four identity
matrix. The physical CBF-QP minimizes deviation from
MPC subject to instantaneous physical barriers with gain $6$ h$^{-1}$
and the worst queue-arrival residual \cite{ames2017}. Its service deadline
is an evaluation criterion. Solver details and validation records are in
\reproductiondetails{the primary experiments}.

\paragraph{Future activation}
\label{supp:detail:sec:delayed-handover-results}
\label{app:sequential-handover}
The seventeen handover contracts comprise the eight extremes
$(d,\bar w_{\rm veh},\eta,p)\in\{1,50\}\times\{60,300\}\times
\{2750,2950\}\times\{4\}$, primary $(1,300,2850,4)$ contract and eight
held-out contracts. The active plan supplies the activation box
$\mathcal I_{t_a}^+=\widehat x_{t_a}+[-d_{t_a},d_{t_a}]$, where
$(d_{t_a})_i=\sum_j|(L^{-1})_{ij}|(r_{t_a})_j$, by outward conversion of its
correlated endpoint bound. No new observation is used. Timed preparation
includes propagation, optimization, checking and repair. A late or
mismatched result retains the existing schedule through its original expiry.

\paragraph{Connected rings}
\label{app:network-transfer}
The ring model has $R=2$ or $3$ regions, classes $c_{ij}$ and one
queue per region. Regional and gate indices start at zero. The regional
total is $N_i=\sum_jc_{ij}$. Local completion is $c_{ii}v(N_i)$, nonlocal
transfer is $u_i c_{ij}v(N_i)$, and admission is $a_i=\min\{60z_i,120g_i\}$. Nonlocal
vehicles move one edge and retain their destination. Half the admitted
traffic is local, and the rest is equally divided among other destinations.
Initial regional/queue totals are 3000/200 vehicles, with half the regional
traffic local. Per region, internal demand is 6000 local and 500 nonlocal
veh/h. Exterior arrivals are 4000 veh/h. Queue 0 has a 30\%
increase from minute 15 to 45. Ring/admission gates start at 0.65/0.50.
After minute 30, ring gates 0/1 become 0.80/0.55 and admission gate 0
becomes 0.60, with other gates unchanged. Class/queue half-widths are
$d/R$ and $\min(d,10)$ vehicles. The regional deadline is 2700 vehicles
at one hour. Physical capacities match the primary model. The plan is
fixed before certification. \reproductiondetails{the ring experiment}
gives the remaining uncertainty pairs, checker settings and all sixteen
outcomes.

\paragraph{Agent histories and scoring}
\label{sec:supp-execution}
The four histories cover direct instruction, restoration after one or
nested interruptions, and interruption expiry at minute 20. The pilot's
presaved cases are Q000 (certified MPC), D013 (repair), and D013 with
replacement tools unavailable (retention). Actual planning starts from
D027, with one-vehicle sensing half-width, 300~veh/h residual bound,
2000-vehicle storage and the 2850-vehicle Region-1 limit at minute 60.
Objectives set weights on $(n_1,n_2,z)$ to $(1,1,1)$, $(4,1,1)$ or
$(1,4,1)$ for network, Region~1 or Region~2. Objective correctness uses
the active operator instruction. Pilot tool correctness uses the saved
checker ordering: certified economic MPC, certified scenario MPC,
checked repair, then retention. Detailed histories, schemas and raw
calls are in \reproductiondetails{the agent protocol}.

\bibliographystyle{elsarticle-num}
\setlength{\bibsep}{0pt}
\bibliography{FOUNDATION_REFERENCES}
\end{document}
